\subsection{External hilbertian sum of hilbertian spaces}

PROPOSITION 1. --- Let \((\mathrm{E}_i)_{i\in \mathbf{I}}\) be a family of hilbertian spaces, \(\mathbf{P}\) the product vector space \(\prod_{i\in \mathbf{I}}\mathrm{E}_i\) and \(\mathbf{E}\) the subset of \(\mathbf{P}\) consisting of all families \(\mathbf{x} = (x_i)_{i\in \mathbf{I}}\) such that \(\sum_{i\in \mathbf{I}}\| x_i\|^2\) is finite.

\begin{enumerate}
\def\labelenumi{\alph{enumi})}
\item
  E is a vector subspace of \(\mathbf{P}\) .
\item
  For every \(\mathbf{x} = (x_i)_{i\in \mathbf{I}}\) and \(\mathbf{y} = (y_i)_{i\in \mathbf{I}}\) in E, the family \((\langle x_i|y_i\rangle)_{i\in \mathbf{I}}\) is summable. If we put \(\langle \mathbf{x}|\mathbf{y}\rangle = \sum_{i\in \mathbf{I}}\langle x_i|y_i\rangle\) , we define a positive separating hermitian form on E.
\item
  For the scalar product so defined, \(\mathbf{E}\) is a hilbertian space; the direct sum \(\mathbf{S}\) of the \(\mathbf{E}_i\) is dense in \(\mathbf{E}\) .
\end{enumerate}

For \(\mathbf{x} = (x_{i})_{i\in \mathbf{I}}\) and \(\mathbf{y} = (y_i)_{i\in \mathbf{I}}\) in E, we have

\[
\left\| x _ {i} + y _ {i} \right\| ^ {2} \leqslant 2 \left(\left\| x _ {i} \right\| ^ {2} + \left\| y _ {i} \right\| ^ {2}\right),
\]

hence \(\mathbf{x} + \mathbf{y} = (x_i + y_i)_{i\in I}\) belongs to E. This proves \(a)\) .

By the Cauchy-Schwarz inequality, we have

\[
\left| \left\langle x _ {i} \mid y _ {i} \right\rangle \right| \leqslant \| x _ {i} \| \cdot \| y _ {i} \| \leqslant \frac {1}{2} \left(\left\| x _ {i} \right\| ^ {2} + \left\| y _ {i} \right\| ^ {2}\right)
\]

hence \(\sum_{i\in \mathbf{I}}|\langle x_i|y_i\rangle | < +\infty\) . If \(x\neq 0\) , we have \(\langle \mathbf{x}|\mathbf{x}\rangle = \sum_{i\in \mathbf{I}}\| x_i\|^2 >0\) , hence assertion \(b)\) follows.

We recall that S is the subspace of P consisting of all families \(\mathbf{x} = (x_{i})_{i \in \mathbf{I}}\) such that the set of all \(i \in I\) for which \(x_{i} \neq 0\) is finite. It follows immediately that S is dense in E; hence it remains to prove that E is complete for the topology \(T_{1}\) obtained by the norm \(\|x\| = \langle x|x\rangle^{1/2}\) . Let \(T_{2}\) be the topology induced on E by the product topology on \(\prod_{i \in I} E_{i}\) . For every r \textgreater{} 0, let \(B_{r}\) be the set of all \(x \in E\) such that \(\|x\| \leqslant r\) . This relation implies that we have \(\sum_{i \in J} \|x_{i}\|^{2} \leqslant r^{2}\) for every finite subset J of I, and so \(B_{r}\) is a closed subset of \(\prod_{i \in I} E_{i}\) , hence also complete. The fact that E is complete for \(T_{1}\) now follows from GT, III, §3, No.~5, cor. 2 to prop. 10.

DEFINITION 1. --- Let \((\mathrm{E}_{i})_{i\in\mathrm{I}}\) be a family of hilbertian spaces. The hilbertian space E defined in prop. 1 is called the external hilbertian sum of the family \((\mathrm{E}_{i})_{i\in\mathrm{I}}\) and written as \(\bigoplus_{i\in I} E_{i}\) or \(\bigoplus_{i\in I} E_{i}^{1}\) .

\(^{1}\) Care must be taken not to confuse this notation with that of the « algebraic » direct sum of the spaces \(E_{i}\) (A, II, § 1, No.~6).

Let \(f_{i}\) be the mapping from \(E_{i}\) into \(E\) which transforms \(z \in E_{i}\) into an element \((x_{k}) \in E\) such that \(x_{k} = 0\) for all \(k \neq i\) and \(x_{i} = z\) ; it is clear that \(f_{i}\) is an isomorphism from the hilbertian space \(E_{i}\) onto a closed vector subspace of \(E\) . We say that \(f_{i}\) is the canonical mapping from \(E_{i}\) into \(E\) and we shall generally identify \(E_{i}\) with its image in \(E\) by this isomorphism. With this convention, \(E_{i}\) and \(E_{k}\) are orthogonal in \(E\) for \(i \neq k\) , and \(E\) is the closed vector subspace generated by the union of the subspaces \(E_{i}\) .

When I is finite, E is the direct sum of the \(E_{i}\) ; since the canonical projector from E onto \(E_{i}\) is continuous for all \(i \in I\) , E is also the topological direct sum of the \(E_{i}\) (GT, III, § 6, No.~2, prop. 2). If \(I = \left(1, n\right)\) , we also write \(E_{1} \oplus E_{2} \oplus \ldots \oplus E_{n}\) instead of \(\bigoplus_{i=1}^{n} E_{i}\) .

Example. --- Let E be a hilbertian space and I a set of indices. Let \(\ell_{\mathrm{E}}^{2}(\mathrm{I})\) denote the external hilbertian sum of the family \((\mathrm{E}_{i})_{i\in\mathrm{I}}\) where \(E_{i}=E\) for all \(i\in I\) . In other words, \(\ell_{\mathrm{E}}^{2}(\mathrm{I})\) is the space of all families \(\mathbf{x}=(x_{i})_{i\in\mathrm{I}}\) of elements of E such that \(\sum_{i\in I}\|x_{i}\|^{2}<+\infty\) , endowed with the scalar product \(\langle x|y\rangle=\sum_{i\in I}\langle x_{i}|y_{i}\rangle\) (space of square summable families of elements of E indexed by I). We put \(\ell^{2}(\mathrm{I})=\ell_{\mathrm{K}}^{2}(\mathrm{I})\) .

\subsection{Hilbertian sum of orthogonal subspaces of a hilbertian space}

DEFINITION 2. --- A hilbertian space E is said to be a hilbertian sum of a family \((\mathbf{E}_{i})_{i\in I}\) of closed vector subspaces of E when :

\begin{enumerate}
\def\labelenumi{\arabic{enumi})}
\item
  for two distinct indices \(i\) , \(k\) in I, the subspaces \(\mathrm{E}_i\) and \(\mathrm{E}_k\) are orthogonal in E;
\item
  the closed vector subspace generated by the union of the \(\mathbf{E}_i\) is E.
\end{enumerate}

THEOREM 1. --- Let E be a hilbertian space which is a hilbertian sum of a family \((\mathrm{E}_{i})_{i\in\mathbf{I}}\) of closed vector subspaces of E. There exists an isomorphism f and only one, from E onto the external hilbertian sum \(\bigoplus_{i\in I}E_{i}=F\) of the family \((\mathrm{E}_{i})\) such that, for all \(i\in I\) , the restriction of \(f\) to \(\mathbf{E}\) is the canonical mapping \(f_{i}\) from \(\mathbf{E}_i\) into \(\mathbf{F}\) .

Let \(S \subset F\) be the « algebraic » direct sum of the \(E_i\) , and let \(g\) be the linear mapping \((x_i)_{i \in I} \mapsto \sum_{i \in I} x_i\) from \(S\) into \(E\) . We shall show that \(g\) is an isomorphism from the prehilbertian space S onto the (prehilbertian) subspace \(g(S)\) of E, generated by the union of the \(E_{i}\) : for, for two elements \(\mathbf{x} = (x_{i})_{i \in \mathrm{I}}\) , \(\mathbf{y} = (y_{i})_{i \in \mathrm{I}}\) , we have

\[
\langle g (\mathbf {x}) | g (\mathbf {y}) \rangle = \left\langle \sum_ {i \in \mathrm{I}} x _ {i} \right| \sum_ {i \in \mathrm{I}} y _ {i} \rangle = \sum_ {(i, k) \in \mathrm{I} \times \mathrm{I}} \left\langle x _ {i} \mid y _ {k} \right\rangle .
\]

But if \(i \neq k\) , \(\langle x_i | y_k \rangle = 0\) by hypothesis, hence

\[
\langle g (\mathbf {x}) | g (\mathbf {y}) \rangle = \sum_ {i \in \mathbf {I}} \langle x _ {i} | y _ {i} \rangle = \langle \mathbf {x} | \mathbf {y} \rangle ;
\]

this proves our assertion. Since S is dense in F and \(g(S)\) dense in E, the isomorphism g extends to an isomorphism \(\overline{g}\) from F onto E (V, p.~8, cor.). It is clear that the inverse isomorphism f of \(\overline{g}\) is the required mapping; its uniqueness follows from the fact that the closed subspace of E generated by the union of the \(E_{i}\) is E itself.

When E is the hilbertian sum of a family \((\mathrm{E}_{i})_{i\in I}\) of subspaces, we shall often identify E with the external hilbertian sum F of the \(E_{i}\) by means of the isomorphism f.~If the set I is finite, saying that E is the hilbertian sum of the family \((\mathrm{E}_{i})_{i\in I}\) means that the \(E_{i}\) are two by two orthogonal and that the vector space E is the direct sum of the family \((\mathrm{E}_{i})_{i\in I}\) of subspaces.

COROLLARY 1. --- Let E be a hilbertian space, which is a hilbertian sum of a family \((\mathbf{E}_{i})_{i\in\mathbf{I}}\) of closed vector subspaces of E; for all \(i\in I\) , let \(p_{E_{i}}\) be the orthoprojector (V, p.~13) from E onto \(E_{i}\) .

\begin{enumerate}
\def\labelenumi{\alph{enumi})}
\tightlist
\item
  For all \(x \in \mathbb{E}\) , the family \((\| p_{\mathbf{E}_i}(x)\|^{2})_{i\in \mathbf{I}}\) is summable in \(\mathbf{R}\) , the family \((p_{\mathbf{E}_i}(x))_{i\in \mathbf{I}}\) is summable in \(\mathbb{E}\) , and we have
\end{enumerate}

\[
\left\| x \right\| ^ {2} = \sum_ {i \in \mathbf {I}} \left\| p _ {\mathrm{E} _ {i}} (x) \right\| ^ {2}, \quad x = \sum_ {i \in \mathbf {I}} p _ {\mathrm{E} _ {i}} (x).
\]

\begin{enumerate}
\def\labelenumi{\alph{enumi})}
\setcounter{enumi}{1}
\item
  Conversely, if \((x_{i})_{i\in \mathbf{I}}\) is a family of elements of \(\mathbf{E}\) such that \(x_{i}\in E_{i}\) for all \(i\in \mathbf{I}\) and \(\sum_{i\in \mathbf{I}}\| x_i\|^2 < + \infty\) , this family is summable, and the sum \(x\) is the only point of \(\mathbf{E}\) for which \(p_{\mathrm{E}_i}(x) = x_i\) for all \(i\in \mathbf{I}\) .
\item
  For every pair of points \(x, y\) of \(\mathbf{E}\) , we have
\end{enumerate}

\[
\langle x | y \rangle = \sum_ {i \in \mathbf {I}} \left\langle p _ {\mathrm{E} _ {i}} (x) \mid p _ {\mathrm{E} _ {i}} (y) \right\rangle .
\]

These properties are in fact obvious for the external hilbertian sum of the \(\mathbf{E}_i\) , and can be transferred to \(\mathbf{E}\) by isomorphism.

COROLLARY 2. --- Let E be a Hausdorff prehilbertian space, \((\mathrm{E}_{i})_{i\in\mathrm{I}}\) a family of complete vector subspaces of E such that, for every pair of distinct indices i, k in I, the subspaces \(E_{i}\) and \(E_{k}\) are orthogonal. Let V be the closed vector subspace of E generated by the union of the \(E_{i}\) . For every \(i\in I\) , let \(p_{E_{i}}\) be the orthoprojector from E onto \(E_{i}\) . Let \(x\in E\) .

\begin{enumerate}
\def\labelenumi{\arabic{enumi})}
\item
  We have \(\sum\left\|p_{E_{i}}(x)\right\|^{2}\leqslant\left\|x\right\|^{2}\) .
\item
  The following conditions are equivalent: a) \(x \in \mathbf{V}\) ; b) \(\sum_{i \in I} \| p_{\mathrm{E}_i}(x) \|^2 = \| x \|^2\) ; c) the family \((p_{\mathrm{E}_i}(x))_{i \in I}\) is summable in E, and we have \(x = \sum_{i \in I} p_{\mathrm{E}_i}(x)\) .
\item
  Suppose V is complete. Then the family \((p_{\mathrm{E}_i}(x))_{i\in \mathbf{I}}\) is summable in E, and
\end{enumerate}

\[
p _ {\mathbf {V}} (x) = \sum_ {i \in \mathbf {I}} p _ {\mathrm{E} _ {i}} (x), \quad \left\| p _ {\mathbf {V}} (x) \right\| ^ {2} = \sum_ {i \in \mathbf {I}} \left\| p _ {\mathrm{E} _ {i}} (x) \right\| ^ {2},
\]

where \(p_{V}\) denotes the orthoprojector from E onto V.

Let \(\hat{\mathbf{E}}\) be the hilbertian space completion of \(\mathbf{E}\) ; we identify \(\mathbf{E}\) with a dense subspace of \(\hat{\mathbf{E}}\) ; the \(\mathbf{E}_i\) , being complete, are closed subspaces of \(\hat{\mathbf{E}}\) . The closure \(\overline{\mathbf{V}}\) of \(\mathbf{V}\) in \(\hat{\mathbf{E}}\) is the closed vector subspace of \(\hat{\mathbf{E}}\) generated by the union of the \(\mathbf{E}_i\) , and \(\mathbf{V} = \overline{\mathbf{V}} \cap \mathbf{E}\) . The space \(\hat{\mathbf{E}}\) is the hilbertian sum of the \(\mathbf{E}_i\) and of the subspace \(\mathbf{W}\) , the orthogonal complement of \(\overline{\mathbf{V}}\) in \(\hat{\mathbf{E}}\) ; put \(x_0 = p_{\mathrm{W}}(x)\) and \(x_i = p_{\mathrm{E}_i}(x)\) for all \(i \in \mathbf{I}\) . By cor. 1, we have \(\| x\|^2 = \| x_0\|^2 + \sum_{i \in \mathbf{I}} \| x_i\|^2\) , and \(x = x_0 + \sum_{i \in \mathbf{I}} x_i\) in \(\hat{\mathbf{E}}\) . This implies assertion 1), and the fact that conditions \(b)\) and \(c)\) of 2) are equivalent to the condition \(x_0 = 0\) , hence to the condition \(x \in \mathbf{V}\) . Finally, if \(\mathbf{V}\) is complete, and if we put \(x' = p_{\mathrm{V}}(x)\) , we have \(x' - x_i = (x - x_i) - (x - p_{\mathrm{V}}(x))\) , hence \(x' - x_i\) is orthogonal to \(\mathbf{E}_i\) , and so \(x_i = p_{\mathrm{E}_i}(x')\) for all \(i \in \mathbf{I}\) ; it is now enough to apply property 2) to the vector \(x'\) .

Remark. --- Let E be a Hausdorff prehilbertian space, \((\mathrm{V}_{i})_{i\in\mathrm{I}}\) a family of vector subspaces of E such that for every pair of distinct indices i, k, the subspaces \(V_{i}\) and \(V_{k}\) are orthogonal. Then, for every \(k\in I\) , the intersection of \(V_{k}\) and of the closed vector subspace \(W_{k}\) generated by the union of the \(V_{i}\) for all \(i\neq k\) reduces to 0 for, if x belongs to \(V_{k}\) and also to \(W_{k}\) , then it is orthogonal to all the \(V_{i}\) for \(i\neq k\) , hence to \(W_{k}\) . In particular, x is orthogonal to itself, hence is zero.

PROPOSITION 2. --- Let E be a hilbertian space and \((\mathrm{V}_{\lambda})_{\lambda\in\mathrm{L}}\) a family of closed vector subspaces of E; for every \(\lambda\in L\) , let \((\mathrm{W}_{\lambda\mu})_{\mu\in\mathrm{M}_{\lambda}}\) be a family of closed vector subspaces of \(V_{\lambda}\) such that \(V_{\lambda}\) is the closed vector subspace generated by the union of this family. In order that E is the hilbertian sum of the family \((\mathrm{W}_{\lambda\mu})_{\lambda\in\mathrm{L},\mu\in\mathrm{M}_{\lambda}}\) , it is necessary and sufficient that E is the hilbertian sum of the family \((\mathrm{V}_{\lambda})_{\lambda\in\mathrm{L}}\) and that, for each \(\lambda\in L\) , \(V_{\lambda}\) is the hilbertian sum of the family \((\mathrm{W}_{\lambda\mu})_{\mu\in\mathrm{M}_{\lambda}}\) (« associativity of the hilbertian sum »). To show that the condition is necessary, it is enough to see that \(V_{\alpha}\) and \(V_{\beta}\) are orthogonal if \(\alpha\neq\beta\) . But, every element of \(W_{\alpha\mu}\) ( \(\mu\in M_{\alpha}\) ) is orthogonal to all the \(W_{\beta\nu}\) ( \(\nu\in M_{\beta}\) ), hence to the closed vector subspace \(V_{\beta}\) which they generate; the same argument then shows that every element of \(V_{\beta}\) is orthogonal to \(V_{\alpha}\) , being orthogonal to all the \(W_{\alpha\mu}\) ( \(\mu\in M_{\alpha}\) ).

To show that the condition is sufficient, it is enough to verify that, if it is satisfied, \(\mathrm{E}\) is equal to the closed vector subspace \(\mathrm{F}\) generated by the union of the \(\mathbf{W}_{\lambda \mu}\) (\(\lambda \in \mathbf{L}\), \(\mu\in M_{\lambda}\)); but, for each \(\lambda\in L\) , F contains the closed vector subspace generated by the union of the \(W_{\lambda\mu}\) such that \(\mu\in M_{\lambda}\) , that is, F contains \(V_{\lambda}\) ; hence F is the closed vector subspace generated by the union of the \(V_{\lambda}\) , which is E by hypothesis.

\subsection{Orthonormal families}

DEFINITION 3. --- In a prehilbertian space, a family \((e_{i})_{i\in\mathbf{I}}\) of vectors is said to be orthogonal if \(e_{i}\) and \(e_{k}\) are orthogonals for all \(i\neq k\) , and is said to be orthonormal, if in addition \(\|e_{i}\|=1\) for all \(i\in I\) .

A subset S of E such that the family defined by the identity mapping from S onto itself is orthonormal is said to be an orthonormal set. If \((e_{i})_{i\in I}\) is an orthonormal family, the mapping \(i\mapsto e_{i}\) is injective; we can then talk indifferently of an orthonormal family or an orthonormal set.

If \((e_i)_{i\in I}\) is an orthonormal family, the complete one dimensional vector subspaces \(\mathbf{D}_i = \mathbf{K}e_i\) are two by two orthogonal. For every \(x\in E\) , the orthogonal projection of \(x\) on \(\mathbf{D}_i\) is \(\lambda_i e_i\) with \(\langle e_i|x - \lambda_i e_i\rangle = 0\) , which gives \(\langle e_i|x\rangle = \lambda_i\langle e_i|e_i\rangle = \lambda_i\) . The results of No.~2 applied to the subspaces \(\mathbf{D}_i\) imply the following propositions:

PROPOSITION 3. --- In a Hausdorff prehilbertian space E, every orthonormal family is topologically independent.

We note that this property immediately follows from the characterization of topologically independent families (IV, p.~1 and II, p.~43, cor. 2), on account of the identification of the dual of E with the completion of E or with the space conjugate to E according as K is equal to R or C (V, p.~17, Remark).

PROPOSITION 4. --- Let E be a Hausdorff prehilbertian space, \((e_{i})_{i\in I}\) an orthonormal family in E, V the closed vector subspace of E generated by the \(e_{i}\) .

\begin{enumerate}
\def\labelenumi{\arabic{enumi})}
\tightlist
\item
  For every \(x \in E\) , we have
\end{enumerate}

\[
\sum_ {i \in \mathbf {I}} \left| \langle e _ {i} | x \rangle \right| ^ {2} \leqslant \| x \| ^ {2}\tag{1}
\]

(Bessel's inequality); here the set of all \(i \in \mathbf{I}\) such that \(\langle e_i|x\rangle \neq 0\) is countable. Moreover, the following conditions are equivalent: a) \(x \in \mathbf{V}\) ; b) \(\|x\|^2 = \sum_{i \in \mathbf{I}} |\langle e_i|x\rangle|^2\) ; c) the family \(\langle e_i|x\rangle\) . \(e_i\) is summable in E, and \(x = \sum_{i \in \mathbf{I}} \langle e_i|x\rangle\) . \(e_i\) .

\begin{enumerate}
\def\labelenumi{\arabic{enumi})}
\setcounter{enumi}{1}
\item
  If \(\mathbf{V}\) is complete, then the family of all \(\langle e_i|x\rangle\) . \(e_i\) is summable in \(\mathbf{E}\) for all \(x\in\mathbf{E}\) , and \(\sum_{i\in\mathbf{I}}\langle e_i|x\rangle\) . \(e_i=p_{\mathbf{V}}(x)\) , \(\sum_{i\in\mathbf{I}}|\langle e_i|x\rangle|^2=\|p_{\mathbf{V}}(x)\|^2\) .
\item
  Suppose V is complete. For every family \((\lambda_{i})_{i\in\mathbf{I}}\) of scalars such that \(\sum_{i\in I}|\lambda_{i}|^{2}<+\infty\) , there exists a unique point \(x\in V\) such that \(\langle e_{i}|x\rangle=\lambda_{i}\) for all \(i\in I\) . If \((\mu_{i})_{i\in I}\) is a second family of scalars such that \(\sum_{i\in I}\left|\mu_{i}\right|^{2}<+\infty\) , and if \(y\in V\) is such that \(\langle e_{i}|y\rangle=\mu_{i}\) for all \(i\in I\) , then \(\langle x|y\rangle=\sum_{i\in I}\overline{\lambda}_{i}\mu_{i}\) .
\end{enumerate}

PROPOSITION 5. --- Let \((e_i)_{i \in I}\) be an orthonormal family in a Hausdorff prehilbertian space E. The following properties are equivalent :

\begin{enumerate}
\def\labelenumi{\alph{enumi})}
\item
  the family \((e_i)\) is total;
\item
  for every \(x \in E\) , the family \(\langle e_{i}|x\rangle\) . \(e_{i}\) is summable in E, and we have \(x = \sum_{i \in I} \langle e_{i}|x\rangle\) . \(e_{i}\) ;
\item
  for every \(x \in E\) ,
\end{enumerate}

\[
\left\| x \right\| ^ {2} = \sum_ {i \in \mathbf {I}} \left| \langle e _ {i} | x \rangle \right| ^ {2}\tag{2}
\]

(Parseval's relation).

When E is hilbertian these conditions are also equivalent to :

\begin{enumerate}
\def\labelenumi{\alph{enumi})}
\setcounter{enumi}{3}
\tightlist
\item
  the relations \(\langle e_i|x\rangle = 0\) for all \(i\in \mathrm{I}\) imply that \(x = 0\) .
\end{enumerate}

The equivalence of conditions a), b), c) follows immediately from prop. 4. When E is hilbertian, the equivalence of conditions a) and d) follows from cor. 1 of V, p.~16.

DEFINITION 4. --- An orthonormal and total family in a Hausdorff prehilbertian space E is called an orthonormal basis of E.

An orthonormal basis of a Hausdorff prehilbertian space E is also an orthonormal basis of the completion of E.

Let \((e_{i})_{i\in I}\) be an orthonormal basis of E; for every \(x\in E\) , the numbers \(\langle e_{i}|x\rangle\) are called, by abuse of language, the coordinates of x with respect to the basis \((e_{i})\) . For every x and y in E, we have

\[
\langle x | y \rangle = \sum_ {i \in \mathbf {I}} \overline {{\langle e _ {i} | x \rangle}} \langle e _ {i} | y \rangle .\tag{3}
\]

An orthonormal basis of E is not, in general, a basis of E over the field K in the sense defined in A, II, p.~25; to avoid any confusion we shall always say that a basis of a prehilbertian space E in the sense of loc. cit. is an algebraic basis of E over K.

Let E and F be two Hausdorff prehilbertian spaces and u a continuous linear mapping from E into F. Let \((e_{i})_{i\in I}\) (resp. \((f_{j})_{j\in J}\) ) be an orthonormal basis of E (resp. F). Put

\[
u _ {j i} = \langle f _ {j} | u (e _ {i}) \rangle
\]

for \(i \in \mathbf{I}\) , \(j \in \mathbf{J}\) . The family \((u_{ji})_{(i,j) \in \mathbf{I} \times \mathbf{J}}\) is called the matrix of \(u\) with respect to the ortho-normal bases \((e_i)\) and \((f_j)\) . Let \(x \in \mathbf{E}\) and \(y = u(x)\) ; if we write \(\xi_i = \langle e_i | x \rangle\) and \(\eta_j = \langle f_j | y \rangle\) for the coordinates of \(x\) and \(y\) respectively, we get \(\eta_j = \sum_{i \in \mathbf{I}} u_{ji} \xi_i\) for all \(j \in \mathbf{J}\) . When \((e_i)\) is an algebraic basis of \(E\) and \((f_j)\) an algebraic basis of \(F\) , our definition is consistent with that of A, II, § 10, No.~4.

Example. --- Let E be the space of all complex valued continuous functions on R, such that \(f(x + n) = f(x)\) for \(x \in R\) and \(n \in Z\) . We assign to E the scalar product defined by

\[
\langle f | g \rangle = \int_ {0} ^ {1} \overline {{{{f (t)}}}} g (t) d t.
\]

Then E is a Hausdorff prehilbertian space, but is not complete. For every integer \(n \in Z\) , let \(e_{n}(x) = \mathbf{e}(nx)\) . It is immediate that the family \((e_{n})_{n \in \mathbf{Z}}\) is orthonormal in E. Moreover, the topology of uniform convergence on E is finer than the topology deduced from the norm \(\|f\|_{2} = \langle f | f \rangle^{1/2}\) . The family \((e_{n})_{n \in \mathbf{Z}}\) is total in E for the uniform convergence (GT, X, § 4, No.~4), and a fortiori in the prehilbertian space E. Hence \((e_{n})_{n \in \mathbf{Z}}\) is an orthonormal basis of E.

\subsection{Orthonormalisation}

THEOREM 2. --- For every orthonormal set L in a hilbertian space E, there exists an orthonormal basis B of E containing L.

In fact, let D be the family of all orthonormal subsets of E, linearly ordered by inclusion; it is immediate that this family has finite character (S, III, § 4, No.~5). Hence there exists a maximal family B in D containing L, by th. 1 of S, III, § 4, No.~5. It remains to prove that B is a total set. If not, there will exist a vector \(y \neq 0\) which is orthogonal to all the vectors of B (V, p.~22, prop. 5), and multiplying y by a suitable scalar, we can assume that \(\|y\| = 1\) ; then, \(B \cup \{y\}\) will be an orthonormal set distinct from B and containing B; this contradicts the definition of B; hence the theorem.

COROLLARY 1. --- In every hilbertian space, there exists an orthonormal basis.

It is enough to apply th. 2 to the case L = ∅.

COROLLARY 2. --- Every hilbertian space is isomorphic to a space \(\ell^2 (\mathbf{I})\) .

More precisely, let \((e_{i})_{i\in I}\) be an orthonormal basis of a hilbertian space E. By props. 4 (V, p.~21) and 5 (V, p.~22), the mapping \(\phi\) defined by

\[
\phi (x) = \big (\langle e _ {i} | x \rangle \big) _ {i \in \mathrm{I}}\tag{4}
\]

is a hilbertian space isomorphism from E onto \(\ell^2 (\mathrm{I})\) . The inverse isomorphism \(\psi\) is defined by

\[
\psi ((\lambda_ {i}) _ {i \in \mathbf {I}}) = \sum_ {i \in \mathbf {I}} \lambda_ {i} e _ {i}.\tag{5}
\]

PROPOSITION 6. --- Let E be a Hausdorff prehilbertian space, and let \((a_{n})_{n\in\mathbb{I}}\) (I an interval of N with origin 1) be a countable (finite or not) independent family of vectors of E. There exists an orthonormal family \((e_{n})_{n\in\mathbb{I}}\) , and only one, in E, with the following properties :

\begin{enumerate}
\def\labelenumi{\arabic{enumi})}
\item
  for every integer \(p \in \mathbf{I}\) , the vector subspace of \(\mathbf{E}\) generated by \(e_1, e_2, \ldots, e_p\) is identical with the vector subspace of \(\mathbf{E}\) generated by \(a_1, a_2, \ldots, a_p\) ;
\item
  for every integer \(p \in \mathbf{I}\) , the number \(\langle a_p | e_p \rangle\) is real and \(> 0\) .
\end{enumerate}

In fact, let \(V_{n}\) be the subspace (of dimension n) generated by \(a_{1}, a_{2}, ..., a_{n}\) . If \(n + 1 \in I\) and \(b_{n+1} = a_{n+1} - p_{\mathrm{V}_{n}}(a_{n+1})\) (where \(p_{V_{n}}\) is the orthoprojector onto the complete subspace \(V_{n}\) ), the line \(Kb_{n+1}\) is the orthogonal of \(V_{n}\) in \(V_{n+1}\) . If the \(e_{n}\) satisfy condition 1) of the proposition, we must have \(e_{n+1} = \lambda b_{n+1}\) ; the condition \(\|e_{n+1}\| = 1\) then implies \(|\lambda|^2 \|b_{n+1}\|^2 = 1\) and the condition \(\langle a_{n+1}|e_{n+1} \rangle > 0\) implies \(\lambda \langle a_{n+1}|b_{n+1} \rangle > 0\) ; this completely determines \(\lambda\) , and we have proved that we can determine, by induction, an orthonormal family \((e_n)_{n\in I}\) , and only one, so as to satisfy conditions 1) and 2) of the proposition.

The sequence \((e_{n})_{n\in\mathbb{I}}\) is said to be obtained by orthonormalisation from the independent family \((a_{n})_{n\in\mathbb{I}}\) . It is clear that the vector subspace generated by the family \((e_{n})\) is identical with the vector subspace generated by the family \((a_{n})\) . In particular, if \((a_{n})\) is a total sequence, so is \((e_{n})\) , which is then an orthonormal basis of E; hence we get:

COROLLARY. --- In every Hausdorff prehilbertian space E satisfying the first axiom of countability, there exists a countable orthonormal basis.

If E satisfies the first axiom of countability, then there exists a total sequence in E, and we can always extract an independent total family from such a sequence (A, II, § 7, No.~1, th. 2).

We can give examples of Hausdorff prehilbertian spaces not having any orthonormal basis (V, p.~70, exerc. 2).

Example. --- Let I be the interval \((-1,1)\) of R and E the vector space of real valued continuous functions on I. Let x denote the canonical injection from I into R, considered as an element of E. By the Stone-Weierstrass theorem the sequence \((x^{n})_{n\in\mathbf{N}}\) is total in E for the topology of uniform convergence GT, X, § 4, No.~2).

Consider E as a real Hausdorff prehilbertian space in which the scalar product is given by

\[
\langle f | g \rangle = \int_ {- 1} ^ {1} f (t) g (t) d t.
\]

The sequence \((x^{n})_{n\in \mathbb{N}}\) is then total in the prehilbertian space E. Let \((\Pi_n)_{n\in \mathbb{N}}\) be the sequence obtained by the orthonormalisation of the sequence \((x^n)_{n\in \mathbb{N}}\) . We can show that \(\Pi_n = (n + \frac{1}{2})^{1 / 2}\mathbf{P}_n\) , where the Legendre polynomial \(\mathbf{P}_n\) is defined by

\[
\mathrm{P} _ {n} (x) = \frac {1}{2 ^ {n} n !} \left(\frac {d}{d x}\right) ^ {n} (x ^ {2} - 1) ^ {n}.
\]

PROPOSITION 7. --- In a hilbertian space E, two orthonormal bases are equipotent.

Let B and C be two orthonormal bases of E. The case when one of the two sets B, C is finite is trivial, since a finite orthonormal basis is an algebraic basis of the space. Suppose therefore that B and C are infinite. For every \(x \in B\) , let \(C_{x}\) be the subset of C consisting of all \(y \in C\) such that \(\langle x | y \rangle \neq 0\) . The set \(C_{x}\) is countable (V, p.~21, prop. 4). For every \(y \in C\) , there exists \(x \in B\) such that \(y \in C_{x}\) , since B is an orthonormal basis and \(y \neq 0\) ; in other words C is the union of the countable sets \(C_{x}\) as x ranges over B. The cardinality of C is hence less than that of N × B, hence less than that of B (S, III, §6, No.~3, cor. 4); similarly, the cardinality of B is less than that of C; this completes the proof.

The cardinality of an arbitrary orthonormal basis of a hilbertian space E is called the hilbertian dimension of E.

COROLLARY 1. --- Given two orthonormal bases in a hilbertian space E, there exists an automorphism of E transforming the first basis into the second.

COROLLARY 2. --- In order that the hilbertian spaces \(\ell^2 (\mathbf{I})\) and \(\ell^2 (\mathbf{J})\) be isomorphic, it is necessary and sufficient that I and J are equipotent.

\section{TENSOR PRODUCT OF HILBERTIAN SPACES}

\subsection{Tensor product of prehilbertian spaces}

Let \(E_{1}\) and \(E_{2}\) be two prehilbertian spaces and let \(F = E_{1} \otimes E_{2}\) be the tensor product of the vector spaces \(E_{1}\) and \(E_{2}\) . Let \(x_{1} \in E_{1}\) and \(x_{2} \in E_{2}\) ; since the mapping \((y_{1}, y_{2}) \mapsto \langle x_{1}|y_{1} \rangle \langle x_{2}|y_{2} \rangle\) from \(E_{1} \times E_{2}\) into K is bilinear, there exists a linear form \(\phi_{x_{1},x_{2}}\) on \(E_{1} \otimes E_{2}\) such that

\[
\phi_ {x _ {1}, x _ {2}} (y _ {1} \otimes y _ {2}) = \langle x _ {1} | y _ {1} \rangle \langle x _ {2} | y _ {2} \rangle\tag{1}
\]

for \(y_{1} \in \mathrm{E}_{1}\) and \(y_{2} \in \mathrm{E}_{2}\) . Let \(z \in \mathrm{F}\) . The mapping \((x_{1}, x_{2}) \mapsto \overline{\phi_{x_{1},x_{2}}(z)}\) from \(\mathrm{E}_1 \times \mathrm{E}_2\) into \(\mathbf{K}\) is bilinear; this can be seen by writing \(z\) in the form \(z = \sum_{i=1}^{n} y_{i,1} \otimes y_{i,2}\) with \(y_{i,1} \in \mathrm{E}_1\) and \(y_{i,2} \in \mathrm{E}_2\) for \(1 \leqslant i \leqslant n\) . Then there exists a linear form \(\psi_z\) on \(\mathrm{F} = \mathrm{E}_1 \otimes \mathrm{E}_2\) such that

\[
\psi_ {z} (x _ {1} \otimes x _ {2}) = \overline {{\phi_ {x _ {1} , x _ {2}} (z)}} \quad (x _ {1} \in \mathrm{E} _ {1}, x _ {2} \in \mathrm{E} _ {2}).\tag{2}
\]

We put \(\Phi(z, t) = \psi_z(t)\) for \(z, t\) in F. We see immediately that \(\Phi\) is a sesquilinear form on \(\mathbf{E}_1 \otimes \mathbf{E}_2\) characterized by

\[
\Phi (x _ {1} \otimes x _ {2}, y _ {1} \otimes y _ {2}) = \langle x _ {1} | y _ {1} \rangle \langle x _ {2} | y _ {2} \rangle\tag{3}
\]

(cf.~A, IX, § 1, No.~11).

PROPOSITION 1. --- The sesquilinear form \(\Phi\) on \(\mathrm{E}_1 \otimes \mathrm{E}_2\) is hermitian and positive, hence assigns the structure of a prehilbertian space to \(\mathrm{E}_1 \otimes \mathrm{E}_2\) . This space is Hausdorff if \(\mathrm{E}_1\) and \(\mathrm{E}_2\) are Hausdorff.

The formula \(\Phi(z, t) = \overline{\Phi(t, z)}\) follows from (3) when \(z = x_1 \otimes x_2\) and \(t = y_1 \otimes y_2\) . The general case is obtained by linearity, hence \(\Phi\) is hermitian.

Suppose \(\mathbf{E}_1\) and \(\mathbf{E}_2\) are Hausdorff; we shall prove that the hermitian form \(\Phi\) is positive and separating. Let \(z = \sum_{i=1}^{n} x_i \otimes y_i\) be a non-zero element of \(\mathrm{F} = \mathrm{E}_1 \otimes \mathrm{E}_2\) . Let \((e_1, ..., e_n)\) be an orthonormal basis of the subspace of \(\mathbf{E}_1\) generated by \(x_1, ..., x_n\) (V, p.~23, cor. 1). There exist elements \(f_1, ..., f_m\) in \(\mathbf{E}_2\) , not all null, such that \(z = \sum_{i=1}^{n} e_i \otimes f_i\) , hence

\[
\begin{array}{l} \Phi (z, z) = \sum_ {i, j = 1} ^ {m} \Phi (e _ {i} \otimes f _ {i}, e _ {j} \otimes f _ {j}) \\ = \sum_ {i, j} \langle e _ {i} | e _ {j} \rangle \langle f _ {i} | f _ {j} \rangle = \sum_ {i = 1} ^ {m} \| f _ {i} \| ^ {2} > 0. \end{array}
\]

For the general case, we shall now prove that \(\Phi\) is positive. Let \(\tilde{E}_{i}\) be the Hausdorff prehilbertian space associated with \(E_{i}\) and let \(\pi_{i}\) be the canonical mapping from \(E_{i}\) onto \(\tilde{E}_{i} (i = 1, 2)\) . Put \(\pi = \pi_{1} \otimes \pi_{2}\) . Let \(\tilde{\Phi}\) be the hermitian form on \(\tilde{E}_{1} \otimes \tilde{E}_{2}\) constructed in the same way as \(\Phi\) . We have

\[
\Phi (z, t) = \Phi (\pi (z), \pi (t)) \quad (z \in \mathrm{F}, t \in \mathrm{F}),
\]

and since \(\tilde{\Phi}\) is positive, so is \(\Phi\) .

Q.E.D.

The prehilbertian space defined in prop. 1 is called the tensor product of the prehilbertian spaces \(E_{1}\) and \(E_{2}\) and is written as \(E_{1} \otimes E_{2}\) . Henceforth we shall write \(\langle z|t\rangle\) for \(\Phi(z,t)\) , and therefore by definition

\[
\langle x _ {1} \otimes x _ {2} | y _ {1} \otimes y _ {2} \rangle = \langle x _ {1} | y _ {1} \rangle \langle x _ {2} | y _ {2} \rangle ;\tag{4}
\]

we shall also write \(\| z\| _2\) or \(\| z\|\) for \(\langle z|z\rangle^{1 / 2}\) . From (4), we get

\[
\left\| x _ {1} \otimes x _ {2} \right\| _ {2} = \left\| x _ {1} \right\|. \left\| x _ {2} \right\|,\tag{5}
\]

then the bilinear mapping \((x_{1}, x_{2}) \mapsto x_{1} \otimes x_{2}\) from \(\mathrm{E}_1 \times \mathrm{E}_2\) into \(\mathrm{E}_1 \otimes \mathrm{E}_2\) is continuous.

For i = 1, 2 let \(F_{i}\) be a vector subspace of \(E_{i}\) , endowed with the induced prehilbertian structure. Then \(F_{1} \otimes F_{2}\) can be identified with a vector subspace of \(E_{1} \otimes E_{2}\) (A, II, § 7, No.~7). Formula (4) shows that \(F_{1} \otimes F_{2}\) with the prehilbertian space structure induced by that of \(E_{1} \otimes_{2} E_{2}\) , is exactly \(F_{1} \otimes_{2} F_{2}\) . We shall henceforth identify \(F_{1} \otimes_{2} F_{2}\) with a prehilbertian subspace of \(E_{1} \otimes_{2} E_{2}\) .

PROPOSITION 2. --- For \(i = 1, 2\) , let \(\mathrm{E}_{i}\) and \(\mathrm{F}_{i}\) be two Hausdorff prehilbertian spaces and let \(u_{i} \in \mathcal{L}(\mathrm{E}_{i}; \mathrm{F}_{i})\) . The linear mapping \(u_{1} \otimes u_{2}\) from \(\mathrm{E}_{1} \otimes_{2} \mathrm{E}_{2}\) into \(\mathrm{F}_{1} \otimes_{2} \mathrm{F}_{2}\) is continuous and we have

\[
\left\| u _ {1} \otimes u _ {2} \right\| = \left\| u _ {1} \right\|. \left\| u _ {2} \right\|.
\]

Consider the positive hermitian form on \(\mathbf{E}_1\) given by

\[
f \left(x _ {1}, y _ {1}\right) = \| u _ {1} \| ^ {2} \left\langle x _ {1} \mid y _ {1} \right\rangle - \left\langle u _ {1} \left(x _ {1}\right) \mid u _ {1} \left(y _ {1}\right) \right\rangle .
\]

By prop. 1 (V, p.~25), there exists a positive hermitian form \(\Phi\) on \(E_1 \otimes E_2\) such that

\[
\begin{array}{r l} \Phi (x _ {1} \otimes x _ {2}, y _ {1} \otimes y _ {2}) & = f (x _ {1}, y _ {1}) \langle x _ {2} | y _ {2} \rangle = \\ & = \| u _ {1} \| ^ {2} \langle x _ {1} \otimes x _ {2} | y _ {1} \otimes y _ {2} \rangle - \langle (u _ {1} \otimes 1) (x _ {1} \otimes x _ {2}) | (u _ {1} \otimes 1) (y _ {1} \otimes y _ {2}) \rangle \end{array}
\]

for \(x_{1}, y_{1}\) in \(E_{1}\) and \(x_{2}, y_{2}\) in \(E_{2}\) . By linearity we have

\[
\Phi (z, t) = \| u _ {1} \| ^ {2} \langle z | t \rangle - \langle (u _ {1} \otimes 1) (z) | (u _ {1} \otimes 1) (t) \rangle
\]

for \(z, t\) in \(\mathrm{E}_1 \otimes \mathrm{E}_2\) . Since \(\Phi\) is positive, we get \(\Phi(z, z) \geqslant 0\) , that is, \(\| (u_1 \otimes 1).z \|_2\) \(\leqslant \| u_1\| \cdot \| z\| _2\) for \(z\in \mathrm{E}_1\otimes_2\mathrm{E}_2\) , or \(\| u_1\otimes 1\| \leqslant \| u_1\|\) . Similarly we prove the inequality \(\| 1\otimes u_2\| \leqslant \| u_2\|\) , and since \(u_{1}\otimes u_{2} = (u_{1}\otimes 1)\circ (1\otimes u_{2})\) , we get

\[
\left\| u _ {1} \otimes u _ {2} \right\| \leqslant \left\| u _ {1} \right\|. \left\| u _ {2} \right\|.
\]

On the other hand,

\[
\begin{array}{l} \| u _ {1} \| \cdot \| u _ {2} \| = \sup _ {\| x _ {1} \| \leqslant 1, \| x _ {2} \| \leqslant 1} \| u _ {1} (x _ {1}) \| \cdot \| u _ {2} (x _ {2}) \| \\ = \sup _ {\| x _ {1} \| \leqslant 1, \| x _ {2} \| \leqslant 1} \| (u _ {1} \otimes u _ {2}) (x _ {1} \otimes x _ {2}) \| _ {2} \leqslant \| u _ {1} \otimes u _ {2} \|. \end{array}
\]

This completes the proof of prop. 2.

Q.E.D.

Let \(E_1, \ldots, E_n\) be prehilbertian spaces ( \(n \geqslant 2\) ). We define the tensor product \(E_1 \otimes_2 \ldots \otimes_2 E_n\) (also denoted by \(\bigotimes_{i=1}^{n} E_i\) ) by induction, by

\[
\mathrm{E} _ {1} \otimes_ {2} \dots \otimes_ {2} \mathrm{E} _ {n} = (\mathrm{E} _ {1} \otimes_ {2} \dots \otimes_ {2} \mathrm{E} _ {n - 1}) \otimes_ {2} \mathrm{E} _ {n}.
\]

Hence, by the definition of scalar product, we have

\[
\left\langle x _ {1} \otimes \dots \otimes x _ {n} \mid y _ {1} \otimes \dots \otimes y _ {n} \right\rangle = \prod_ {i = 1} ^ {n} \left\langle x _ {i} \mid y _ {i} \right\rangle ,\tag{6}
\]

and in particular \(^{1}\)

\[
\left\| x _ {1} \otimes \dots \otimes x _ {n} \right\| _ {2} = \left\| x _ {1} \right\| \dots \left\| x _ {n} \right\|,\tag{7}
\]

for \(x_{i},y_{i}\) in \(\mathrm{E}_i\) ( \(1\leqslant i\leqslant n\) ). If the \(\mathrm{E}_i\) are Hausdorff, then so is \(\mathrm{E}_1\otimes_2\dots \otimes_2\mathrm{E}_n\) .

Let \(F_{1}, \ldots, F_{n}\) be prehilbertian spaces and \(u_{i} \in \mathcal{L}(E_{i}; F_{i})\) for \(1 \leqslant i \leqslant n\) . By induction on n, prop. 2 implies that \(u_{1} \otimes \ldots \otimes u_{n}\) is a continuous linear mapping from \(E_{1} \otimes_{2} \ldots \otimes_{2} E_{n}\) into \(F_{1} \otimes_{2} \ldots \otimes_{2} F_{n}\) and that

\[
\left\| u _ {1} \otimes \dots \otimes u _ {n} \right\| = \left\| u _ {1} \right\| \dots \left\| u _ {n} \right\|.\tag{8}
\]

Let \(\sigma \in \mathfrak{S}_n\) be a permutation of the set \(\{1, 2, \dots, n\}\) . Because of (6), the linear mapping \(p_{\sigma}\) from \(\mathrm{E}_1 \otimes_2 \dots \otimes_2 \mathrm{E}_n\) onto \(\mathrm{E}_{\sigma^{-1}(1)} \otimes_2 \dots \otimes_2 \mathrm{E}_{\sigma^{-1}(n)}\) characterized by

\[
p _ {\sigma} \left(x _ {1} \otimes \dots \otimes x _ {n}\right) = x _ {\sigma^ {- 1} (1)} \otimes \dots \otimes x _ {\sigma^ {- 1} (n)}\tag{9}
\]

is a prehilbertian space isomorphism (« commutativity of tensor product »).

Similarly, consider a partition of \(\{1, 2, \dots , n\}\) into \(m\) consecutive intervals \(\mathrm{I}_{1}, \dots , \mathrm{I}_{m}\) with \(\mathrm{I}_{k} = [a_{k}, a_{k + 1} - 1]\) for \(1 \leqslant k \leqslant m\). Put

\[
\mathrm{F} _ {k} = \bigotimes_ {i = a _ {k}} ^ {a _ {k + 1} - 1} \mathrm{E} _ {i} (1 \leqslant k \leqslant m).
\]

\(^{1}\) Here again we put \(\|z\|_{2} = \langle z|z\rangle^{1/2}\) for z in \(E_{1} \otimes_{2} \ldots \otimes_{2} E_{n}\) .

The canonical isomorphism from \(F_{1} \otimes \ldots \otimes F_{m}\) onto \(E_{1} \otimes \ldots \otimes E_{n}\) which transforms \(\bigotimes_{k=1}^{m} \bigotimes_{i=a_{k}}^{a_{k+1}-1} x_{i}\) into \(x_{1} \otimes \ldots \otimes x_{n}\) (A, II, §3, No.~9) is a prehilbertian space isomorphism (« associativity of the tensor product »).

\subsection{Hilbertian tensor product of hilbertian spaces}

DEFINITION 1. --- Let \(E_{1}, \ldots, E_{n}\) be hilbertian spaces. The completion of the Hausdorff prehilbertian space \(E_{1} \otimes_{2} \ldots \otimes_{2} E_{n}\) is called the hilbertian tensor product of the \(E_{i}\) and is denoted by \(E_{1} \hat{\otimes}_{2} \ldots \hat{\otimes}_{2} E_{n}\) (or \(\hat{\otimes}_{2} E_{i}\) ).

Let \(F_{1}, \ldots, F_{n}\) be hilbertian spaces and \(u_{i} \in \mathcal{L}(E_{i}, F_{i})\) for \(1 \leqslant i \leqslant n\) . The continuous linear mapping \(u_{1} \otimes \ldots \otimes u_{n}\) then extends to a continuous linear mapping \(u_{1} \hat{\otimes}_{2} \ldots \hat{\otimes}_{2} u_{n}\) from \(E_{1} \hat{\otimes}_{2} \ldots \hat{\otimes}_{2} E_{n}\) into \(F_{1} \hat{\otimes}_{2} \ldots \hat{\otimes}_{2} F_{n}\) . We have

\[
\| u _ {1} \hat {\otimes} _ {2} \dots \hat {\otimes} _ {2} u _ {n} \| = \| u _ {1} \| \dots \| u _ {n} \|\tag{10}
\]

by formula (8) of V, p.~27. Moreover, if \(1_{E}\) denotes the identity mapping of any hilbertian space E, we have

\[
1 _ {\mathrm{E} _ {1}} \hat {\otimes} _ {2} \dots \hat {\otimes} _ {2} 1 _ {\mathrm{E} _ {n}} = 1 _ {\mathrm{E}} \quad \text { with } \quad \mathrm{E} = \mathrm{E} _ {1} \hat {\otimes} _ {2} \dots \hat {\otimes} _ {2} \mathrm{E} _ {n}.\tag{11}
\]

Finally, if \(\mathbf{G}_1, \ldots, \mathbf{G}_n\) are hilbertian spaces and \(v_i \in \mathcal{L}(\mathbf{F}_i; \mathbf{G}_i)\) for \(1 \leqslant i \leqslant n\) , we get

\[
\left(v _ {1} \circ u _ {1}\right) \hat {\otimes} _ {2} \dots \hat {\otimes} _ {2} \left(v _ {n} \circ u _ {n}\right) = \left(v _ {1} \hat {\otimes} _ {2} \dots \hat {\otimes} _ {2} v _ {n}\right) \circ \left(u _ {1} \hat {\otimes} _ {2} \dots \hat {\otimes} _ {2} u _ {n}\right).\tag{12}
\]

We leave to the reader the task of formulating the «commutativity» and the «associativity» of the hilbertian tensor product, in analogy with what has been said above for prehilbertian spaces.

Remark. --- Let \(E_{1}, \ldots, E_{n}\) be Hausdorff prehilbertian spaces, and \(\hat{E}_{1}, \ldots, \hat{E}_{n}\) their respective completions. Then \(E_{1} \otimes_{2} \ldots \otimes_{2} E_{n}\) is a prehilbertian subspace of \(\hat{E}_{1} \otimes_{2} \ldots \otimes_{2} \hat{E}_{n}\) . Since the mapping \((x_{1}, \ldots, x_{n}) \mapsto x_{1} \otimes \ldots \otimes x_{n}\) from \(\hat{E}_{1} \times \cdots \times \hat{E}_{n}\) into \(\hat{E}_{1} \otimes_{2} \ldots \otimes_{2} \hat{E}_{n}\) is continuous, \(E_{1} \otimes_{2} \ldots \otimes_{2} E_{n}\) is dense in \(\hat{E}_{1} \otimes_{2} \ldots \otimes_{2} \hat{E}_{n}\) . A fortiori the completion of \(E_{1} \otimes_{2} \ldots \otimes_{2} E_{n}\) is precisely the hilbertian space \(\hat{E}_{1} \otimes_{2} \ldots \otimes_{2} \hat{E}_{n}\) . This completion is sometimes simply written as \(E_{1} \otimes_{2} \ldots \otimes_{2} E_{n}\) (or \(\bigotimes_{1 \leq i \leq n} E_{i}\) ).

PROPOSITION 3. --- Let \(E_{1}, \ldots, E_{n}\) be hilbertian spaces. Suppose that for \(1 \leqslant i \leqslant n\) the space \(E_{i}\) is a hilbertian sum of a family \((\mathrm{E}_{i,\alpha})_{\alpha \in \mathbf{A}(i)}\) of closed vector subspaces. Then \(E_{1} \otimes_{2} \ldots \otimes_{2} E_{n}\) is a hilbertian sum of the family of subspaces \(E_{1,\alpha_{1}} \hat{\otimes}_{2} \ldots \hat{\otimes}_{2} E_{n,\alpha_{n}}\) with \((\alpha_{1}, \ldots, \alpha_{n})\) ranging over \(A(1) \times \cdots \times A(n)\) .

By formula (6) of V, p.~27, the subspaces \(\mathrm{E}_{1,\alpha_1}\hat{\otimes}_2\ldots \hat{\otimes}_2\mathrm{E}_{n,\alpha_n}\) of \(\mathrm{E}_1\hat{\otimes}_2\ldots \hat{\otimes}_2\mathrm{E}_n\) are mutually orthogonal. For every integer \(i\) between 1 and \(n\) , the set \(\bigcup_{\alpha \in \mathbf{A}(i)}\mathrm{E}_{i,\alpha}\) is total in \(\mathrm{E}_i\) , and the multilinear mapping \((x_{1},\dots,x_{n})\mapsto x_{1}\otimes \dots \otimes x_{n}\) is continuous. It follows that the union of the subspaces \(\mathrm{E}_{1,\alpha_1}\hat{\otimes}_2\ldots \hat{\otimes}_2\mathrm{E}_{n,\alpha_n}\) is total, hence prop. 3.

COROLLARY 1. --- For \(1 \leqslant i \leqslant n\) , let \((e_{i,\alpha})_{\alpha \in \mathbf{A}(i)}\) be an orthonormal basis of \(\mathbf{E}_i\) . Then the family of vectors \(e_{1,\alpha_1} \otimes \ldots \otimes e_{n,\alpha_n}\) as \((\alpha_1, \ldots, \alpha_n)\) ranges over \(\mathbf{A}(1) \times \cdots \times \mathbf{A}(n)\) is an orthonormal basis of \(\mathbf{E}_1 \hat{\otimes}_2 \ldots \hat{\otimes}_2 \mathbf{E}_n\) .

COROLLARY 2. --- Let \(E_{1}\) and \(E_{2}\) be two hilbertian spaces, and \((e_{i})_{i\in I}\) an orthonormal basis of \(E_{1}\) . Let \((y_{i})_{i\in I}\) be a family of elements of \(E_{2}\) , such that \(\sum_{i\in I}\|y_{i}\|^{2}<+\infty\) . Then the family \((e_{i}\otimes y_{i})_{i\in I}\) is summable in \(E_{1}\hat{\otimes}_{2}E_{2}\) ; moreover, every element of \(E_{1}\hat{\otimes}_{2}E_{2}\) can be written uniquely in the form \(\sum_{i\in I}e_{i}\otimes y_{i}\) with \(\sum_{i\in I}\|y_{i}\|^{2}<+\infty\) .

Let \(F_{i}\) be the line in \(E_{1}\) generated by the \(e_{i} (i \in I)\) . Then \(E_{1}\) is the hilbertian sum of the family of subspaces \((F_{i})_{i \in I}\) . By prop. 3, the space \(E_{1} \hat{\otimes} E_{2}\) is the hilbertian sum of the family of subspaces \((F_{i} \hat{\otimes}_{2} E_{2})_{i \in I}\) , hence cor. 2 follows.

Examples. --- 1) By cor. 1, the space \(\ell^2 (\mathrm{I})\hat{\otimes}_2\ell^2 (\mathrm{J})\) is canonically isomorphic to \(\ell^2 (\mathrm{I}\times \mathrm{J})\) , the tensor product \(\mathbf{x}\otimes \mathbf{y}\) of \(\mathbf{x} = (x_i)_{i\in \mathrm{I}}\) and \(\mathbf{y} = (y_j)_{j\in \mathrm{J}}\) can be identified with the family \((x_iy_j)_{i\in \mathrm{I},j\in \mathrm{J}}\) . Similarly, by cor. 2, \(\ell^2 (\mathrm{I})\hat{\otimes}_2\mathrm{E}\) can be identified with \(\ell_{\mathrm{E}}^{2}(\mathrm{I})\) , in such a way that we have \((x_{i})_{i\in \mathrm{I}}\otimes y = (x_{i}y)_{i\in \mathrm{I}}\) for every \(\mathbf{y}\) in the hilbertian space E.

* 2) Let X be a Hausdorff topological space, and \(\mu\) a positive measure on X. Let E be a hilbertian space. We can identify \(L^2(X, \mu) \otimes_2 E\) with \(L_E^2(X, \mu)\) in a canonical way: if \(f\) is the class of the square integrable scalar function \(f\) on X, and if \(a\) belongs to E, then \(f \otimes a\) is the class of the function \(x \mapsto f(x).a\) with values in E.

Let Y be a Hausdorff topological space and v a positive measure on Y. In an analogous manner, we can identify the hilbertian spaces \(\mathrm{L}^{2}(\mathrm{X},\mu)\hat{\otimes}_{2}\mathrm{L}^{2}(\mathrm{Y},v)\) and \(\mathrm{L}^{2}(\mathrm{X}\times\mathrm{Y},\mu\otimes v)\) ; then \(\dot{f}\otimes\dot{g}\) can be identified with the class of the function \((x,y)\mapsto f(x)g(y)\) on \(X\times Y\) .

\subsection{Symmetric hilbertian powers}

Let E be a hilbertian space, and let n be a positive integer. Let \(\hat{\mathbf{T}}^{n}(\mathrm{E})\) or \(E^{\otimes_{n}}\) denote the tensor product of n hilbertian spaces each equal to E. In other words, \(\hat{\mathbf{T}}^{n}(\mathrm{E})\) is the completion of the space \(\mathbf{T}^{n}(\mathrm{E}) = \mathrm{E} \otimes \ldots \otimes \mathrm{E}\) (n factors) for the Hausdorff prehilbertian space structure defined by

\[
\left\langle x _ {1} \otimes \dots \otimes x _ {n} \mid y _ {1} \otimes \dots \otimes y _ {n} \right\rangle = \prod_ {i = 1} ^ {n} \left\langle x _ {i} \mid y _ {i} \right\rangle .
\]

If \((e_i)_{i\in I}\) is an orthonormal basis of \(E\) , the family of vectors \(e_{i_1} \otimes \ldots \otimes e_{i_n}\) for \(i_1, \ldots, i_n\) in \(I\) , is an orthonormal basis of \(\hat{\mathbf{T}}^n(E)\) (V, p.~29, cor. 1). We get \(\hat{\mathbf{T}}^\circ(E) = K\) .

Let \(\sigma \in \mathfrak{S}_n\) be a permutation of the set \(\{1, 2, ..., n\}\) . By V, p.~27, there exists an automorphism \(p_{\sigma}\) of \(\hat{\mathbf{T}}^n(\mathbf{E})\) characterized by

\[
p _ {\sigma} (x _ {1} \otimes \dots \otimes x _ {n}) = x _ {\sigma^ {- 1} (1)} \otimes \dots \otimes x _ {\sigma^ {- 1} (n)}.
\]

We have \(p_{\sigma \tau} = p_{\sigma}p_{\tau}\) for \(\sigma, \tau\) in \(\mathfrak{S}_n\) , and consequently the endomorphism \(\Pi_n = \frac{1}{n!}\sum_{\sigma \in \mathfrak{S}_n}p_\sigma\) of the vector space \(\hat{\mathbf{T}}^n (\mathrm{E})\) is the orthoprojector onto the subspace of all elements left invariant by \(\mathfrak{S}_n\) . However (A, III, § 6, No.~3), \(\Pi_n\) maps the « algebraic » tensor product \(\mathbf{T}^n (\mathrm{E})\) onto the subspace \(\mathbf{T}\mathbf{S}^n (\mathrm{E})\) of all symmetric tensors of order \(n\) . In other words, the image of \(\Pi_n\) is the completion of the space \(\mathbf{T}\mathbf{S}^n (\mathrm{E})\) endowed with a scalar product induced by that of \(\mathbf{T}^n (\mathrm{E})\) ; this completion will be denoted by \(\widehat{\mathbf{T}\mathbf{S}}^n (\mathrm{E})\) .

Let \(\mathbf{S}^n (\mathrm{E})\) be the \(n\) th symmetric power of the vector space E (A, III, § 6, No.~1). The canonical mapping from \(\mathbf{T}^n (\mathrm{E})\) onto \(\mathbf{S}^n (\mathrm{E})\) defines by restriction an isomorphism \(\lambda_{n}\) from \(\mathbf{TS}^n (\mathrm{E})\) onto \(\mathbf{S}^n (\mathrm{E})\) . We verify immediately that the inverse isomorphism is given by

\[
\mu_ {n} (x _ {1} \dots x _ {n}) = \Pi_ {n} (x _ {1} \otimes \dots \otimes x _ {n}) = \frac {1}{n !} \sum_ {\sigma \in \mathfrak {S} _ {n}} x _ {\sigma^ {- 1} (1)} \otimes \dots \otimes x _ {\sigma^ {- 1} (n)}\tag{13}
\]

for \(x_{1},\ldots ,x_{n}\) in E.

We define a Hausdorff prehilbertian space structure on \(\mathbf{S}^{n}(\mathrm{E})\) by putting

\[
\langle u | v \rangle = n! \left\langle \mu_ {n} (u) | \mu_ {n} (v) \right\rangle .\tag{14}
\]

We then have (compare with formula (29) of A, III, § 11, No.~5)

\[
\left\langle x _ {1} \dots x _ {n} \mid y _ {1} \dots y _ {n} \right\rangle = \sum_ {\sigma \in \mathfrak {S} _ {n}} \prod_ {i = 1} ^ {n} \left\langle x _ {i} \mid y _ {\sigma (i)} \right\rangle ,\tag{15}
\]

and in particular

\[
\langle x ^ {n} | y ^ {n} \rangle = n! \langle x | y \rangle^ {n}.\tag{16}
\]

Let \(\hat{\mathbf{S}}^n (\mathrm{E})\) denote the completion of the pre-Hilbertian space \(\mathbf{S}^n (\mathrm{E})\) and \(\hat{\mathbf{S}} (\mathrm{E})\) the external hilbertian sum of the hilbertian spaces \(\hat{\mathbf{S}}^n (\mathrm{E})\) . We can show (V, p.~73, exerc. 1) that the multiplication in the algebra \(\mathbf{S}(\mathrm{E})\) cannot be extended by continuity to \(\mathbf{S}(\mathrm{E})\) , unless E is just 0.

PROPOSITION 4. --- Let \((e_{i})_{i\in I}\) be an orthonormal basis of the hilbertian space E. For every \(\alpha\) in \(\mathbf{N}^{(1)}\) , put

\[
z _ {\alpha} = \prod_ {i \in I} e _ {i} ^ {\alpha_ {i}} / (\alpha_ {i}!) ^ {1 / 2}.\tag{17}
\]

Then \((z_{\alpha})_{\alpha \in \mathbf{N}^{(1)}}\) is an orthonormal basis of \(\hat{\mathbf{S}} (\mathrm{E})\) .

Let \(E_{0}\) be the vector subspace of E generated by the vectors \(e_{i}\) for i ranging over I. Then the \(z_{\alpha}\) form a basis of the vector space \(\mathbf{S}(E_{0})\) (A, III, § 6, No.~6). But \(E_{0}\) is dense in E, and the multilinear mapping \((x_{1},...,x_{n})\mapsto x_{1}\ldots x_{n}\) from \(E\times\cdots\times E\) into \(\mathbf{S}(E)\) is continuous for all \(n\geqslant1\) ; hence \(\mathbf{S}(E_{0})\) is dense in \(\mathbf{S}(E)\) . It is now enough to prove that the family of the \(z_{\alpha}\) is orthonormal. First observe that \(\hat{\mathbf{S}}^{n}(\mathrm{E})\) and \(\hat{\mathbf{S}}^{m}(\mathrm{E})\) are orthogonal for \(n \neq m\) . Hence it is enough to prove the formula

\[
\langle z _ {\alpha} | z _ {\beta} \rangle = \left\{ \begin{array}{l l} 1 & \text {if} \quad \alpha = \beta \\ 0 & \text {if} \quad \alpha \neq \beta \end{array} \right.
\]

when \(|\alpha| = \sum_{i \in I} \alpha_i\) and \(|\beta| = \sum_{i \in I} \beta_i\) are equal to the same integer \(n\) .

Consider a partition \((\mathbf{P}_i)_{i\in I}\) of the set \(\{1,2,\dots,n\}\) such that \(\operatorname{Card} \mathbf{P}_i = \alpha_i\) for all \(i\in I\) . Put \(x_k = e_i\) if \(k\) belongs to \(\mathbf{P}_i\) , then \(x_1\ldots x_n = \prod_{i\in I}e_i^{\alpha_i}\) . Similarly we define \((\mathbf{Q}_i)_{i\in I}\) and \(y_k\) in such a way that \(\operatorname{Card} \mathbf{Q}_i = \beta_i\) and \(y_1\ldots y_n = \prod_{i\in I}e_i^{\beta_i}\) . Since the \(e_i\) are mutually orthogonal, we have \(\langle x_k|y_{\sigma(k)}\rangle = 0\) except if there exists an indice \(i\in I\) such that \(k\in \mathbf{P}_i\) and \(\sigma (k)\in \mathbf{Q}_i\) . By formula (15), we then have \(\langle x_1\ldots x_n|y_1\ldots y_n\rangle = 0\) unless there exists a permutation \(\sigma \in \mathfrak{S}_n\) such that \(\sigma (\mathbf{P}_i) = \mathbf{Q}_i\) for all \(i\in I\) , which implies that \(\alpha = \beta\) . Then \(\langle z_{\alpha}|z_{\beta}\rangle = 0\) for \(\alpha \neq \beta\) . The same argument proves that \(\| x_1\ldots x_n\|^2\) is equal to the number of the \(\sigma \in \mathfrak{S}_n\) such that \(\sigma (\mathbf{P}_i) = \mathbf{P}_i\) for all \(i\in I\) , hence equal to \(\prod_{i=1}^{\infty}\alpha_i!\) . We get \(\| z_\alpha \| = 1\) , and the proposition is proved.

COROLLARY. --- Suppose that the hilbertian space E is the direct sum of the orthogonal subspaces M and N. The canonical isomorphism g from \(\mathbf{S}(\mathbf{M}) \otimes \mathbf{S}(\mathbf{N})\) onto \(\mathbf{S}(\dot{\mathbf{E}})\) (A, III, §6, No.~6) extends uniquely to a hilbertian space isomorphism h from \(\hat{\mathbf{S}}(\mathbf{M}) \otimes_{2} \hat{\mathbf{S}}(\mathbf{N})\) onto \(\hat{\mathbf{S}}(\mathbf{E})\) .

Let \((e_{i})_{i\in I}\) (resp. \((f_{j})_{j\in J}\) ) be an orthonormal basis of the hilbertian space M (resp. N) and let \(M_{0}\) (resp. \(N_{0}\) ) be the vector subspace of E generated by the vectors \(e_{i}\) (resp. \(f_{j}\) ). Put \(E_{0}=M_{0}+N_{0}\) and let \(g_{0}\) be the canonical isomorphism from \(S(M_{0})\otimes S(N_{0})\) onto \(S(E_{0})\) . Put

\[
z _ {\alpha} = \prod_ {i \in \mathbf {I}} e _ {i} ^ {\alpha_ {i}} / (\alpha_ {i}!) ^ {1 / 2}, t _ {\beta} = \prod_ {j \in \mathbf {J}} f _ {j} ^ {\beta _ {j}} / (\beta_ {j}!) ^ {1 / 2}
\]

for \(\alpha \in \mathbf{N}^{(I)}\) and \(\beta \in \mathbf{N}^{(J)}\) . By prop. 4, we have thus defined the orthonormal bases \((z_{\alpha})_{\alpha \in \mathbf{N}^{(I)}}\) for \(\hat{\mathbf{S}}(\mathbf{M})\) , \((t_{\beta})_{\beta \in \mathbf{N}^{(J)}}\) for \(\hat{\mathbf{S}}(\mathbf{N})\) and \((z_{\alpha}t_{\beta})_{\alpha \in \mathbf{N}^{(I)},\beta \in \mathbf{N}^{(J)}}\) for \(\hat{\mathbf{S}}(\mathbf{E})\) . Since we have \(z_{\alpha}t_{\beta} = g_0(z_{\alpha}\otimes t_{\beta})\) , and since the elements \(z_{\alpha}\otimes t_{\beta}\) form an orthonormal basis of \(\hat{\mathbf{S}}(\mathbf{M})\hat{\otimes}_2\hat{\mathbf{S}}(\mathbf{N})\) (V, p.~29, cor. 1), we see that \(g_0\) extends to a hilbertian space isomorphism \(h:\hat{\mathbf{S}}(\mathbf{M})\hat{\otimes}_2\hat{\mathbf{S}}(\mathbf{N})\to \hat{\mathbf{S}}(\mathbf{E})\) . By the construction, we have

\[
h (x _ {1} \dots x _ {m} \otimes y _ {1} \dots y _ {n}) = x _ {1} \dots x _ {m} y _ {1} \dots y _ {n}
\]

for all vectors \(x_{1}, \ldots, x_{m}\) in \(M_{0}\) and \(y_{1}, \ldots, y_{n}\) in \(N_{0}\) . By continuity, the same relation also holds for the vectors \(x_{1}, \ldots, x_{m}\) in M and the vectors \(y_{1}, \ldots, y_{n}\) in N; in other words, h extends g. The uniqueness of h is clear.

Let E and F be two hilbertian spaces and \(u \in \mathcal{L}(E; F)\) . The linear mapping \(\hat{\mathbf{T}}^{n}(u) = u \otimes_{2} \ldots \otimes_{2} u\) (n factors) from \(\hat{\mathbf{T}}^{n}(E)\) into \(\hat{\mathbf{T}}^{n}(F)\) is continuous with norm \(\|u\|^{n}\) (V, p.~28, formula (10)). Moreover, formulas (13) and (14) of V, p.~30, show that there exists an isomorphism \(\phi_{n,\mathrm{E}}\) from \(\hat{\mathbf{S}}^n (\mathrm{E})\) onto the subspace \(\widehat{\mathbf{T}\mathbf{S}}^n (\mathrm{E})\) of \(\hat{\mathbf{T}}^n (\mathrm{E})\) , and only one, such that

\[
\phi_ {n, \mathrm{E}} (x _ {1} \dots x _ {n}) = \frac {1}{(n !) ^ {1 / 2}} \sum_ {\sigma \in \mathfrak {S} _ {n}} x _ {\sigma (1)} \otimes \dots \otimes x _ {\sigma (n)} \quad (x _ {1},..., x _ {n} \text { in } \mathrm{E}).\tag{18}
\]

Hence there exists a continuous linear mapping \(\hat{\mathbf{S}}^{n}(u)\) from \(\hat{\mathbf{S}}^{n}(\mathrm{E})\) into \(\hat{\mathbf{S}}^{n}(\mathrm{F})\) and only one which makes the following diagram commutative

\[
\begin{array}{c c c} \hat {\mathbf {S}} ^ {n} (\mathrm{E}) & \xrightarrow {\phi_ {n , \mathrm{E}}} & \hat {\mathbf {T}} ^ {n} (\mathrm{E}) \\ \hat {\mathbf {S}} ^ {n} (u) \Big \downarrow & & \Big \downarrow \hat {\mathbf {T}} ^ {n} (u) \\ \hat {\mathbf {S}} ^ {n} (\mathrm{F}) & \xrightarrow {\phi_ {n , \mathrm{F}}} & \hat {\mathbf {T}} ^ {n} (\mathrm{F}) \end{array}
\]

We now prove the formula

\[
\left\| \hat {\mathbf {S}} ^ {n} (u) \right\| = \| u \| ^ {n}.\tag{19}
\]

We clearly have \(\| \hat{\mathbf{S}}^n (u)\| \leqslant \| \hat{\mathbf{T}}^n (u)\| = \| u\| ^n\) . Further, for all \(x\in \mathrm{E}\) , we have \(\hat{\mathbf{S}}^n (u)(x^n) = u(x)^n,\| x^n\| = (n!)^{1 / 2}\| x\| ^n\) and \(\| u(x)^n\| = (n!)^{1 / 2}\| u(x)\| ^n\) , which gives

\[
\left\| \hat {\mathbf {S}} ^ {n} (u) \right\| \| x \| ^ {n} \geqslant \left\| u (x) \right\| ^ {n};
\]

it follows immediately that \(\| \hat{\mathbf{S}}^n (u)\| \geqslant \| u\| ^n\) , hence formula (19).

It is clear that we have the formulas

\[
\hat {\mathbf {S}} ^ {n} (1 _ {\mathrm{E}}) = 1 _ {\hat {\mathbf {S}} ^ {n} (\mathrm{E})}\tag{20}
\]

\[
\hat {\mathbf {S}} ^ {n} (v \circ u) = \hat {\mathbf {S}} ^ {n} (v) \circ \hat {\mathbf {S}} ^ {n} (u) \quad \text { for } \quad v \in \mathscr {L} (\mathrm{F}; \mathrm{G}).\tag{21}
\]

Finally, \(\hat{\mathbf{S}}^n (u)\) coincides on \(\mathbf{S}^n (\mathrm{E})\) with the linear mapping \(\mathbf{S}^n (u):\mathbf{S}^n (\mathrm{E})\to \mathbf{S}^n (\mathrm{F})\) defined in A, III, § 6, No.~2 since it transforms \(x_{1}\ldots x_{n}\) into \(u(x_{1})\ldots u(x_{n})\) for every \(x_{1},\ldots ,x_{n}\) in E.

Examples. --- * 1) Let \(d \geqslant 1\) be an integer and \(\omega\) a positive function on \(\mathbf{R}^{d}\) , locally integrable with respect to the Lebesgue measure \(\mu\) . Let E be the hilbertian space \(\mathrm{L}^{2}(\mathbf{R}^{d}, \omega, \mu)\) , and let \(\mathbf{S} = \mathbf{S}(\mathrm{E})\) . Then \(\mathbf{S}\) can be identified with the space of all sequences \(f = (f_{n})_{n \geqslant 0}\) , where each \(f_{n}\) is a function on \((\mathbf{R}^{d})^{n}\) which is measurable with respect to the Lebesgue measure \(\mu \otimes \ldots \otimes \mu\) ( \(n\) factors) and invariant under the permutations of the \(n\) factors in \((\mathbf{R}^{d})^{n}\) , and is such that

\[
\| \mathbf {f} \| ^ {2} = \sum_ {n = 0} ^ {\infty} n! \int_ {\mathbf {R} ^ {d}} \dots \int_ {\mathbf {R} ^ {d}} \left| f _ {n} (\mathbf {x} _ {1}, \dots , \mathbf {x} _ {n}) \right| ^ {2} \omega (\mathbf {x} _ {1}) \dots \omega (\mathbf {x} _ {n}) d \mathbf {x} _ {1} \dots d \mathbf {x} _ {n}\tag{22}
\]

is finite. The norm \(\| \mathbf{f}\|\) in \(\mathbf{S}\) is defined by formula (22). The hilbertian space \(\mathbf{S}\) defined above is called the symmetric Fock space corresponding to the weight \(\omega_{*}\)

* 2) Let X be a Hausdorff topological space, \(\mu\) a positive measure of norm 1 on X and E a hilbertian subspace of the real hilbertian space \(L_{\mathbb{R}}^{2}(X,\mu)\) . We say that E is a gaussian space if the following equivalent conditions are satisfied:

\begin{enumerate}
\def\labelenumi{\alph{enumi})}
\item
  for all \(f \in \mathrm{E}\) , we have \(\int_{\mathbf{x}} e^{if} d\mu = \exp(-\|f\|^2/2)\) ;
\item
  for all \(f \in E\) with norm 1, the image of the measure \(\mu\) under f is the measure
\end{enumerate}

\[
(2 \pi) ^ {- 1 / 2} e ^ {- x ^ {2} / 2} d x.
\]

on \(\mathbf{R}\) .

Suppose E is a gaussian space. Let \(f_1, \ldots, f_n\) be functions whose classes \(f_i\) belong to E. We define a function: \(f_1 \ldots f_n\) : on X (called « Wick's product » of \(f_1, \ldots, f_n\) ) by the formula

\[
: f _ {1} \dots f _ {n} := \sum_ {0 \leqslant 2 p \leqslant n} (- 1) ^ {p} \sum_ {\sigma \in I _ {p}} \prod_ {i = 1} ^ {p} \left\langle f _ {\sigma (2 i - 1)} \mid f _ {\sigma (2 i)} \right\rangle \prod_ {j = 2 p + 1} ^ {n} f _ {\sigma (j)},\tag{23}
\]

where \(\mathbf{I}_p\) is the set of permutations \(\sigma\) of \(\{1,2,\dots,n\}\) such that we have

\[
\begin{array}{l} \sigma (1) <   \sigma (2), \dots , \sigma (2 p - 1) <   \sigma (2 p) \\ \sigma (1) <   \sigma (3) <   \dots <   \sigma (2 p - 1) \\ \sigma (2 p + 1) <   \sigma (2 p + 2) <   \dots <   \sigma (n). \end{array}
\]

Then there exists an isomorphism \(\phi\) from \(\hat{\mathbf{S}}(\mathrm{E})\) onto a hilbertian subspace of \(\mathrm{L}_{\mathbb{R}}^{2}(\mathrm{X},\mu)\) which transforms the product \(\dot{f}_1\dots \dot{f}_n\) of \(\dot{f}_1,\dots ,\dot{f}_n\) , calculated in \(\hat{\mathbf{S}}(\mathrm{E})\) , into \((:f_1\dots f_n:)\). Suppose that X is a Souslin space and that there exists a countable family \((f_{n})\) of functions whose classes belong to E and which separate the points of X. Then \(\phi\) is an isomorphism from \(\hat{\mathbf{S}}(\mathrm{E})\) onto \(\mathrm{L}_{\mathbb{R}}^{2}(\mathrm{X},\mu)\) .

\subsection{Exterior hilbertian powers}

Let E be a hilbertian space and n a positive integer. For every permutation \(\sigma\in\mathfrak{S}_{n}\) , let \(\varepsilon_{\sigma}\) denote its signature; put \(a_{n}=\frac{1}{n!}\sum_{\sigma\in\mathfrak{S}_{n}}\varepsilon_{\sigma}p_{\sigma}\) in \(\mathcal{L}(\hat{\mathbf{T}}^{n}(\mathrm{E}))\) (V, p.~29). It is immediate that \(a_{n}\) is an orthoprojector, whose image \(\overline{\mathbf{A}_{n}^{\prime}(\mathrm{E})}\) is the closure in \(\hat{\mathbf{T}}^{n}(\mathrm{E})\) of the space \(\mathbf{A}_{n}^{\prime}(\mathrm{E})\) of all antisymmetric tensors of order \(n(\mathrm{A},\mathrm{III},\S7,\mathrm{No}.4)\) . There exists an isomorphism \(\pi_{n}\) from \(\Lambda^{n}(\mathrm{E})\) onto \(\mathbf{A}_{n}^{\prime}(\mathrm{E})\) which is characterized by

\[
\pi_ {n} (x _ {1} \wedge \ldots \wedge x _ {n}) = \mathbf {a} _ {n} (x _ {1} \otimes \ldots \otimes x _ {n}) = \frac {1}{n !} \sum_ {\sigma \in \mathfrak {S} _ {n}} \varepsilon_ {\sigma} x _ {\sigma (1)} \otimes \ldots \otimes x _ {\sigma (n)}\tag{24}
\]

for \(x_{1}, \ldots, x_{n}\) in E. We can now define a Hausdorff prehilbertian space structure on \(\Lambda^{n}(E)\) by putting

\[
\langle u | v \rangle = n! \left\langle \pi_ {n} (u) \mid \pi_ {n} (v) \right\rangle .\tag{25}
\]

More explicitly, we have (compare with formula (30) of A, III, § 11, No.~5).

\[
\langle x _ {1} \wedge \dots \wedge x _ {n} | y _ {1} \wedge \dots \wedge y _ {n} \rangle = \det (\langle x _ {i} | y _ {j} ^ {\prime} \rangle)\tag{26}
\]

for all \(x_{1},\ldots ,x_{n}\) and \(y_{1},\ldots ,y_{n}\) in E.

Let \(\hat{\mathbf{A}}^n (\mathrm{E})\) denote the completion of the prehilbertian space \(\mathbf{A}^n (\mathrm{E})\) , and \(\hat{\mathbf{A}} (\mathrm{E})\) the external hilbertian sum of the hilbertian spaces \(\hat{\mathbf{A}}^n (\mathrm{E})\) .

Example. --- * With the notations of example 1 of V, p.~32, we can identify the hilbertian space \(\symbf{\Lambda}(\mathrm{E})\) with the set of all sequences \((f_n)_{n\geqslant 0}\) of measurable functions for which the number \(\| \mathbf{f}\|\) defined in (22) is finite, and where each function \(f_{n}\) is antisymmetric, that is, satisfies the relation

\[
f _ {n} (\mathbf {x} _ {\sigma (1)}, \dots , \mathbf {x} _ {\sigma (n)}) = \varepsilon_ {\sigma} f _ {n} (\mathbf {x} _ {1}, \dots , \mathbf {x} _ {n})
\]

for every permutation \(\sigma \in \mathfrak{S}_n\) . The hilbertian space \(\hat{\Lambda}(\mathrm{E})\) is called the antisymmetric Fock space corresponding to the weight \(\omega\) . *

PROPOSITION 5. --- Let \((e_i)_{i \in I}\) be an orthonormal basis of the hilbertian space E. Endow I with a totally ordered structure. Then the set of all elements \(e_{i_1} \wedge \ldots \wedge e_{i_n}\) for \(i_1 < \cdots < i_n\) is an orthonormal basis of \(\hat{\Lambda}^n(E)\) .

We know (A, III, § 7, No.~8) that the elements in question form a basis of the vector space \(\Lambda^n (\mathrm{E}_0)\) where \(\mathrm{E}_0\) is the vector subspace of E generated by the vectors \(e_i\) . Further, for \(i_1 < \dots < i_n\) , the matrix of scalar products \(\langle e_{i_k}|e_{i_\ell}\rangle\) is the unit matrix of order \(n\) ; by (26), we have \(\| e_{i_1}\wedge \ldots \wedge e_{i_n}\| = 1\) . Finally, if \((i_1,\dots,i_n)\) and \((j_1,\dots,j_n)\) are two distinct, strictly increasing sequences of elements of I, then there exists an element \(j_\ell\) distinct from \(i_1,\ldots ,i_n\) and so we have \(\langle e_{i_k}|e_{j_\ell}\rangle = 0\) for \(1\leqslant k\leqslant n\) , and by (26), \(\langle e_{i_1}\wedge \ldots \wedge e_{i_n}|e_{j_1}\wedge \ldots \wedge e_{j_n}\rangle = 0\) . In other words, the family of elements \(e_{i_1}\wedge \ldots \wedge e_{i_n}\) , for \(i_1 < \dots < i_n\) , is orthonormal.

But \(E_{0}\) is dense in E, and the mapping \((x_{1},...,x_{n})\mapsto x_{1}\wedge\ldots\wedge x_{n}\) from \(E\times\cdots\times E\) into \(\Lambda^{n}(E)\) is continuous. Consequently, \(\Lambda^{n}(E_{0})\) is dense in \(\Lambda^{n}(E)\) , and proposition 5 follows.

COROLLARY. --- Suppose that the hilbertian space E is the direct sum of two orthogonal subspaces M and N. The canonical isomorphism g from \(\symbf{\Lambda}(\mathbf{M})\otimes\symbf{\Lambda}(\mathbf{N})\) onto \(\symbf{\Lambda}(\mathbf{E})\) (A, III, § 7, No.~7) extends in a unique way to a hilbertian space isomorphism from \(\hat{\symbf{\Lambda}}(\mathbf{M})\hat{\otimes}_{2}\hat{\symbf{\Lambda}}(\mathbf{N})\) onto \(\hat{\symbf{\Lambda}}(\mathbf{E})\) .

The proof is analogous to that of the corollary of prop. 4 (V, p.~31).

Let \(\mathbf{E}\) and \(\mathbf{F}\) be two hilbertian spaces and \(u \in \mathcal{L}(\mathbf{E};\mathbf{F})\) . We shall show, as in the case of symmetric powers \(\hat{\mathbf{S}}^n (\mathbf{E})\) (V, p.~32) that the linear mapping \(\symbf{\Lambda}^n (u)\) from \(\symbf{\Lambda}^n (\mathbf{E})\) into \(\symbf{\Lambda}^n (\mathbf{F})(\mathbf{A},\mathbf{III},\S 7,\mathbf{No}.4)\) extends to a continuous linear mapping \(\hat{\symbf{\Lambda}}^n (u)\) from \(\hat{\symbf{\Lambda}}^n (\mathbf{E})\) into \(\hat{\symbf{\Lambda}}^n (\mathbf{F})\) . We have the relations

\[
\hat {\symbf {\Lambda}} ^ {n} (1 _ {\mathrm{E}}) = 1 _ {\hat {\symbf {\Lambda}} ^ {n} (\mathrm{E})},\tag{27}
\]

\[
\hat {\symbf {\Lambda}} ^ {n} (v \circ u) = \hat {\symbf {\Lambda}} ^ {n} (v) \circ \hat {\symbf {\Lambda}} ^ {n} (u) \quad \text { if } v \text { belongs   to } \mathscr {L} (\mathrm{F}; \mathrm{G}),\tag{28}
\]

\[
\left\| \hat {\symbf {\Lambda}} ^ {n} (u) \right\| \leqslant \| u \| ^ {n}.\tag{29}
\]

In general we do not have equality in formula (29) (TS, IV, § 6). Finally, we have an isomorphism \(\psi_{n} = \psi_{n,\mathrm{E}}\) from \(\hat{\mathbf{A}}^{n}(\mathrm{E})\) onto the subspace \(\overline{\mathbf{A}_{n}^{\prime}(\mathrm{E})}\) of \(\hat{\mathbf{T}}^{n}(\mathrm{E})\) defined by

\[
\psi_ {n} (x _ {1} \wedge \dots \wedge x _ {n}) = \frac {1}{(n !) ^ {1 / 2}} \sum_ {\sigma \in \mathfrak {S} _ {n}} \varepsilon_ {\sigma} x _ {\sigma (1)} \otimes \dots \otimes x _ {\sigma (n)}.\tag{30}
\]

\subsection{Exterior multiplication}

Let E be a hilbertian space. For every integer \(n \geqslant 0\) , let \(\theta_{n}\) be the canonical mapping from \(\mathbf{T}^{n}(\mathbf{E})\) onto \(\symbf{\Lambda}^{n}(\mathbf{E})\) ; then

\[
\theta_ {n} (x _ {1} \otimes \ldots \otimes x _ {n}) = x _ {1} \wedge \ldots \wedge x _ {n}\tag{31}
\]

for \(x_{1}, \ldots, x_{n}\) in E. Let \(p\) and \(q\) be two positive integers; on account of formulas (30) and (31) we have

\[
u \wedge v = \theta_ {p + q} \left(\frac {1}{(p !) ^ {1 / 2}} \psi_ {p} (u) \otimes \frac {1}{(q !) ^ {1 / 2}} \psi_ {q} (v)\right)\tag{32}
\]

for \(u \in \Lambda^p(\mathbf{E})\) and \(v \in \Lambda^q(\mathbf{E})\) . Since \(\| \theta_n \| \leqslant (n!)^{1/2}\) , we get the inequality

\[
\| u \wedge v \| \leqslant \left(\frac {(p + q) !}{p ! q !}\right) ^ {1 / 2} \| u \| \cdot \| v \|\tag{33}
\]

for \(u \in \symbf{\Lambda}^p(\mathrm{E})\) and \(v \in \symbf{\Lambda}^q(\mathrm{E})\) . Consequently, the mapping \((u, v) \mapsto u \wedge v\) extends by continuity to a bilinear mapping from \(\hat{\symbf{\Lambda}}^p(\mathrm{E}) \times \hat{\symbf{\Lambda}}^q(\mathrm{E})\) into \(\hat{\symbf{\Lambda}}^{p + q}(\mathrm{E})\) , with a norm at most equal to \(\left(\frac{(p + q)!}{p!q!}\right)^{1/2}\) (cf.~V, p.~73, exerc. 2). We again denote this by \((u, v) \mapsto u \wedge v\) .

PROPOSITION 6. --- Let E be a hilbertian space. We have

\[
\| x \wedge u \| \leqslant \| x \| \cdot \| u \|\tag{34}
\]

for \(x\in \mathbf{E}\) and \(u\in \hat{\Lambda} (\mathrm{E})\)

It is clearly enough to consider the case \(\| x\| = 1\)

Let F be the hilbertian subspace of E consisting of all vectors orthogonal to x. Since E is the hilbertian sum of F and the line K.x, it follows from the corollary of V, p.~34 that the mapping \((v,w)\mapsto v+x\wedge w\) is a hilbertian space isomorphism from \(\hat{\mathbf{A}}(\mathbf{F})\oplus\hat{\mathbf{A}}(\mathbf{F})\) onto \(\hat{\mathbf{A}}(\mathbf{E})\) . If \(u=v+x\wedge w\) with v,w in \(\hat{\mathbf{A}}(\mathbf{F})\) , we have \(x\wedge u=x\wedge v\) , hence \(\|x\wedge u\|=\|v\|\leqslant(\|v\|^{2}+\|w\|^{2})^{1/2}=\|u\|\) .

COROLLARY 1. --- a) Let \(x_1, \ldots, x_n\) be elements of the hilbertian space E. We have

\[
\left\| x _ {1} \wedge \dots \wedge x _ {n} \right\| \leqslant \left\| x _ {1} \right\| \dots \left\| x _ {n} \right\|,\tag{35}
\]

where equality holds only if one of the \(x_{i}\) is null, or if the sequence \((x_{1},\dots,x_{n})\) is orthogonal.

\begin{enumerate}
\def\labelenumi{\alph{enumi})}
\setcounter{enumi}{1}
\tightlist
\item
  Let \(x_{1}, \ldots, x_{n}, y_{1}, \ldots, y_{n}\) be elements of the hilbertian space E. We have
\end{enumerate}

\[
\left| \det (\langle x _ {i}, y _ {j} \rangle) \right| \leqslant \| x _ {1} \| \dots \| x _ {n} \| \cdot \| y _ {1} \| \dots \| y _ {n} \|;\tag{36}
\]

if the vectors \(x_{i}\) and \(y_{i}\) are non-null; equality holds in (36) if and only if \((x_{1},\ldots ,x_{n})\) and \((y_{1},\dots,y_{n})\) are each an orthogonal basis of the same vector subspace of \(\mathbf{E}\) .

The inequality (35) follows from prop. 6 by induction on n; the inequality (36) can be deduced by applying the Cauchy-Schwarz inequality in \(\hat{\Lambda}^{n}(E)\) and the formula (26) of V, p.~33.

Suppose that the sequence \((x_{1},\dots,x_{n})\) is orthogonal. Then

\[
\left\| x _ {1} \wedge \dots \wedge x _ {n} \right\| ^ {2} = \det (\langle x _ {i} | x _ {j} \rangle) = \prod_ {i = 1} ^ {n} \left\| x _ {i} \right\| ^ {2}
\]

since \(\langle x_i|x_j\rangle = 0\) for \(i\neq j\) .

Now suppose that the vectors \(x_{1}, \ldots, x_{n}\) are not null and do not form an orthogonal sequence. Since \(\| x_{1} \wedge \ldots \wedge x_{n} \|\) depends symmetrically on the vectors \(x_{1}, \ldots, x_{n}\) , we can assume that \(x_{1}\) is not orthogonal to the subspace \(F\) of \(E\) generated by \(x_{2}, \ldots, x_{n}\) , and that \(F\) is not simply 0. We can decompose \(x_{1}\) in the form \(x_{1}' + y\) with \(y \neq 0\) in \(F\) and \(x_{1}'\) orthogonal to \(F\) ; then \(\| x_{1}' \| < \| x_{1} \|\) . But

\[
x _ {1} \wedge x _ {2} \wedge \ldots \wedge x _ {n} = x _ {1} ^ {\prime} \wedge x _ {2} \wedge \ldots \wedge x _ {n},
\]

and so

\[
\begin{array}{r l} \left\| x _ {1} \wedge \dots \wedge x _ {n} \right\| & \leqslant \left\| x _ {1} ^ {\prime} \right\| \left\| x _ {2} \right\| \dots \left\| x _ {n} \right\| \\ & <   \left\| x _ {1} \right\| \left\| x _ {2} \right\| \dots \left\| x _ {n} \right\|. \end{array}
\]

Suppose that the vectors \(x_{i}\) and the vectors \(y_{i}\) are not null. The equality in relation (36) is equivalent to the conjunction of the equalities

\[
\left| \left\langle x _ {1} \wedge \dots \wedge x _ {n} \mid y _ {1} \wedge \dots \wedge y _ {n} \right\rangle \right| = \| x _ {1} \wedge \dots \wedge x _ {n} \| \cdot \| y _ {1} \wedge \dots \wedge y _ {n} \|\tag{37}
\]

\[
\left\| x _ {1} \wedge \dots \wedge x _ {n} \right\| = \left\| x _ {1} \right\| \dots \left\| x _ {n} \right\|, \quad \left\| y _ {1} \wedge \dots \wedge y _ {n} \right\| = \left\| y _ {1} \right\| \dots \left\| y _ {n} \right\|.\tag{38}
\]

By the first part of the proof, the equalities (38) imply that each of the sequences \((x_{1},\ldots ,x_{n})\) and \((y_{1},\ldots ,y_{n})\) is orthogonal, which in turn implies that \(x_{1}\wedge \ldots \wedge x_{n}\neq 0\) and that \(y_{1}\wedge \ldots \wedge y_{n}\neq 0\) . Under these conditions, relation (37) implies that there exists a scalar \(\lambda \neq 0\) such that \(y_{1}\wedge \ldots \wedge y_{n} = \lambda x_{1}\wedge \ldots \wedge x_{n}\) (V, p.~3, Remark 1); in other words, that \((x_{1},\dots,x_{n})\) and \((y_{1},\dots,y_{n})\) are bases of the same vector subspace of E (A, III, § 11, No.~13).

COROLLARY 2. --- Let \((a_{ij})_{1\leqslant i,j\leqslant n}\) be a hermitian matrix, with complex elements, and with determinant D. Suppose that the inequality

\[
\sum_ {i, j = 1} ^ {n} a _ {i j} \overline {{z}} _ {i} z _ {j} \geqslant 0\tag{39}
\]

holds for all complex numbers \(z_{1}, \ldots, z_{n}\) . Then

\[
0 \leqslant \mathrm{D} \leqslant a _ {1 1} \dots a _ {n n}.\tag{40}
\]

Suppose \(\mathbf{D}\) is non-null; the equality \(\mathbf{D} = a_{11} \ldots a_{nn}\) holds if and only if \(a_{ij} = 0\) for all \(i \neq j\) .

Let \(\Phi\) be the hermitian form on the vector space \(\mathbf{C}^n\) given by

\[
\Phi (\mathbf {z}, \mathbf {z} ^ {\prime}) = \sum_ {i, j = 1} ^ {n} a _ {i j} \overline {{{{z}}}} _ {i} z _ {j} ^ {\prime}
\]

for \(\mathbf{z} = (z_1, \dots, z_n)\) and \(\mathbf{z}' = (z_1', \dots, z_n')\) in \(\mathbf{C}^n\) . By hypothesis, \(\Phi\) is positive.

First assume that \(\Phi\) is separating, that is, that D is non-null. If \((e_1, \ldots, e_n)\) is the canonical basis of \(\mathbf{C}^n\) , we have \(\Phi(\mathbf{e}_i, \mathbf{e}_j) = a_{ij}\) , and cor. 2 follows immediately from cor. 1, \(a)\) by putting \(x_i = \mathbf{e}_i\) .

Since \(a_{ii} = \Phi(e_i, e_i) \geqslant 0\) , we also have inequality (40) if \(D = 0\) .

COROLLARY 3 (« Hadamard's inequalities »). --- Let \((a_{ij})_{1 \leqslant i,j \leqslant n}\) be a matrix with complex elements, and with determinant D. Put

\[
c _ {i} = \left(\sum_ {j = 1} ^ {n} | a _ {i j} | ^ {2}\right) ^ {1 / 2} \quad \text {for} \quad 1 \leqslant i \leqslant n,
\]

and \(m = \sup_{i,j}|a_{ij}|\) . Then we have

\[
| \mathbf {D} | \leqslant c _ {1} \dots c _ {n} \leqslant m ^ {n}. n ^ {n / 2}.\tag{41}
\]

If \(\mathbf{D} \neq 0\) , in order that \(|\mathbf{D}| = c_1 \ldots c_n\) it is necessary and sufficient that the rows \(y_i = (a_{ij})_{1 \leqslant j \leqslant n}\) of the matrix \((a_{ij})_{1 \leqslant i,j \leqslant n}\) are two by two orthogonal vectors.

Let the space \(C^{n}\) be assigned a scalar product defined by

\[
\langle \mathbf {z} | \mathbf {z} ^ {\prime} \rangle = \sum_ {i = 1} ^ {n} \overline {{z}} _ {i} z _ {i} ^ {\prime}.
\]

Let \((\mathbf{x}_{1},\ldots,\mathbf{x}_{n})\) be the canonical basis of \(C^{n}\) and \(y_{i}\) the vector with components \(a_{ij}\) for \(1\leqslant j\leqslant n\) . We have \(\|x_{i}\|=1\) and \(\|y_{i}\|=c_{i}\) for \(1\leqslant i\leqslant n\) ; also \(\langle x_{i}|y_{j}\rangle=a_{ji}\) . The inequality \(|D|\leqslant c_{1}\ldots c_{n}\) and the condition of equality are then particular cases of V, p.~35, cor. 1. Obviously we have \(c_{i}\leqslant m.n^{1/2}\) , hence \(c_{1}\ldots c_{n}\leqslant m^{n}.n^{n/2}\) .

\section{SOME CLASSES OF OPERATORS IN HILBERTIAN SPACES}

Throughout this paragraph, \(I_{E}\) denotes the identity mapping of a hilbertian space E. The composition \(v \circ u\) of two linear mappings will usually be denoted by vu or v.u.

\subsection{Adjoint}

PROPOSITION 1. --- Let E and F be two hilbertian spaces. For every mapping \(u \in \mathcal{L}(E; F)\) , there exists a unique mapping \(u^{*} \in \mathcal{L}(F; E)\) such that

\[
\langle u (x) | y \rangle_ {\mathrm{F}} = \langle x | u ^ {*} (y) \rangle_ {\mathrm{E}}\tag{1}
\]

for all \(x \in E\) and all \(y \in F\) . The mapping \(u \mapsto u^{*}\) from \(\mathcal{L}(\mathrm{E}, \mathrm{F})\) into \(\mathcal{L}(\mathrm{F}; \mathrm{E})\) is bijective, isometric and semi-linear (with respect to the automorphism \(\xi \mapsto \bar{\xi}\) of K).

Let \(\mathcal{S}(\mathrm{E},\mathrm{F})\) be the space of all continuous sesquilinear forms on \(\mathbf{E}\times \mathbf{F}\) , endowed with the norm

\[
\| \Phi \| = \sup _ {\| x \| \leqslant 1, \| y \| \leqslant 1} | \Phi (x, y) |.\tag{2}
\]

We define the space \(\mathcal{S}(\mathrm{F},\mathrm{E})\) similarly. We defined (V, p.~16, cor. 2) a Banach space isomorphism from \(\mathcal{L}(\mathrm{E};\mathrm{F})\) onto \(\mathcal{S}(\mathrm{F},\mathrm{E})\) , denoted by \(u\mapsto\Phi_{u}\) and characterized by

\[
\Phi_ {u} (y, x) = \langle y | u (x) \rangle_ {\mathrm{F}} (x \in \mathrm{E}, y \in \mathrm{F}).\tag{3}
\]

In an analogous way we define an isomorphism from \(\mathcal{L}(\mathrm{F},\mathrm{E})\) onto \(\mathcal{S}(\mathrm{E},\mathrm{F})\) . Finally we define a mapping \(\Phi \mapsto \Phi^{*}\) from \(\mathcal{S}(\mathrm{F},\mathrm{E})\) onto \(\mathcal{S}(\mathrm{E},\mathrm{F})\) by

\[
\Phi^ {*} (x, y) = \overline {{\Phi (y , x)}} \quad (x \in \mathrm{E}, y \in \mathrm{F}).\tag{4}
\]

This mapping is bijective, semi-linear and isometric. But formula (1) translates as \(\Phi_{u^{*}} = (\Phi_{u})^{*}\) , hence the proposition.

DEFINITION 1. --- Let E and F be two hilbertian spaces. For every continuous linear mapping \(u: \mathbf{E} \to \mathbf{F}\) , the continuous linear mapping from F into E defined by formula (1) is called the adjoint of \(u\) and is denoted by \(u^*\) .

We have

\[
(u + v) ^ {*} = u ^ {*} + v ^ {*}\tag{5}
\]

\[
(\lambda u) ^ {*} = \overline {{{{\lambda}}}} u ^ {*}\tag{6}
\]

\[
(u ^ {*}) ^ {*} = u\tag{7}
\]

\[
(1 _ {\mathrm{E}}) ^ {*} = 1 _ {\mathrm{E}}\tag{8}
\]

\[
(w u) ^ {*} = u ^ {*} w ^ {*};\tag{9}
\]

in all these formulas, u and v belong to \(\mathcal{L}(\mathrm{E};\mathrm{F})\) , \(\lambda\) is in K, and w in \(\mathcal{L}(\mathrm{F};\mathrm{G})\) where G is a hilbertian space. Formulas (5) and (6) mean that \(u\mapsto u^{*}\) is semi-linear. Formula (8) is obvious. To prove (7), we take the conjugate of the two members of (1), which gives \(\langle u^{*}(y)|x\rangle = \langle y|u(x)\rangle\) , and this proves that u is the adjoint of \(u^{*}\) . Finally, with the notations of (9), we have, for all \(z \in G\)

\[
\langle w (u (x)) | z \rangle = \langle u (x) | w ^ {*} (z) \rangle = \langle x | u ^ {*} (w ^ {*} (z)) \rangle ,
\]

hence \(u^{*}w^{*}\) is the adjoint of wu.

Let \(u: \mathrm{E} \to \mathrm{F}\) be a bijective and continuous linear mapping; then it is also bicontinuous (I, p.~19, cor. 1). From (8) and (9) we immediately deduce that \(u^*\) is bijective and bicontinuous and that

\[
(u ^ {- 1}) ^ {*} = (u ^ {*}) ^ {- 1}.\tag{10}
\]

PROPOSITION 2. --- For every \(u \in \mathcal{L}(\mathrm{E};\mathrm{F})\) , we have

\[
\| u ^ {*} u \| = \| u u ^ {*} \| = \| u \| ^ {2} = \| u ^ {*} \| ^ {2}.\tag{11}
\]

By prop. 1, \(\| u^{*}\| = \| u\|\) , hence \(\| u^{*}u\| \leqslant \| u^{*}\| \cdot \| u\| \leqslant \| u\|^2\) . On the other hand,

\[
\| u \| ^ {2} = \sup _ {\| x \| \leqslant 1} \| u (x) \| ^ {2} = \sup _ {\| x \| \leqslant 1} \langle u (x) | u (x) \rangle = \sup _ {\| x \| \leqslant 1} \langle x | u ^ {*} u (x) \rangle \leqslant \| u ^ {*} u \|,
\]

hence \(\| u^{*}u\| = \| u\|^2\) . Replacing \(u\) by \(u^{*}\) , we get \(\| uu^{*}\| = \| u^{*}\|^2\) , hence (11) follows since \(\| u\| = \| u^{*}\|\) .

Let \(E_{1}, \ldots, E_{n}\) and \(F_{1}, \ldots, F_{n}\) be hilbertian spaces, and for every integer i between 1 and n, let \(u_{i}\) be a continuous linear mapping from \(E_{i}\) into \(F_{i}\) . Then

\[
(u _ {1} \hat {\otimes} _ {2} \dots \hat {\otimes} _ {2} u _ {n}) ^ {*} = u _ {1} ^ {*} \hat {\otimes} _ {2} \dots \hat {\otimes} _ {2} u _ {n} ^ {*}.\tag{12}
\]

Let \(v\) be the continuous linear mapping \(u_1 \hat{\otimes}_2 \ldots \hat{\otimes}_2 u_n\) from

\[
\mathbf {E} = \mathbf {E} _ {1} \hat {\otimes} _ {2} \dots \hat {\otimes} _ {2} \mathbf {E} _ {n} \quad \text {into} \quad \mathbf {F} = \mathbf {F} _ {1} \hat {\otimes} _ {2} \dots \hat {\otimes} _ {2} \mathbf {F} _ {n}
\]

and w the continuous linear mapping \(u_{1}^{*}\hat{\otimes}_{2}\ldots\hat{\otimes}_{2}u_{n}^{*}\) from F into E. It is enough to prove the equality \(\langle y|v(x)\rangle=\langle w(y)|x\rangle\) for \(x\in E\) and \(y\in F\) . By linearity and continuity, we reduce to the case when x and y have the following form

\[
x = x _ {1} \otimes \ldots \otimes x _ {n}, y = y _ {1} \otimes \ldots \otimes y _ {n}
\]

with \(x_{i} \in E_{i}\) and \(y_{i} \in F_{i}\) for \(1 \leqslant i \leqslant n\) . From the definition of scalar product in a tensor product (V, p.~27, formula (6)), we then get

\[
\langle y | v (x) \rangle = \prod_ {i = 1} ^ {n} \left\langle y _ {i} \mid u _ {i} \left(x _ {i}\right) \right\rangle = \prod_ {i = 1} ^ {n} \left\langle u _ {i} ^ {*} \left(y _ {i}\right) \mid x _ {i} \right\rangle = \left\langle w (y) \mid x \right\rangle .
\]

This proves our assertion.

Let E and F be two hilbertian spaces, \(u \in \mathcal{L}(\mathrm{E}; \mathrm{F})\) and n a positive integer. If we put \(u_{1} = \cdots = u_{n} = u\) in formula (12) we obtain the result that the continuous linear mapping \(\hat{\mathbf{T}}^{n}(u^{*})\) from \(\hat{\mathbf{T}}^{n}(\mathrm{F})\) into \(\hat{\mathbf{T}}^{n}(\mathrm{E})\) is the adjoint of the continuous linear mapping \(\hat{\mathbf{T}}^{n}(u)\) from \(\hat{\mathbf{T}}^{n}(\mathrm{E})\) into \(\hat{\mathbf{T}}^{n}(\mathrm{F})\) . The formulas

\[
\hat {\mathbf {S}} ^ {n} (u) ^ {*} = \hat {\mathbf {S}} ^ {n} (u ^ {*}), \quad \hat {\symbf {\Lambda}} ^ {n} (u) ^ {*} = \hat {\symbf {\Lambda}} ^ {n} (u ^ {*})\tag{13}
\]

can be established in the same way as formula (12), on account of the definition of the scalar product in \(\hat{\mathbf{S}}^n (\mathrm{E})\) (V, p.~30, formula (15)) and in \(\hat{\Lambda}^n (\mathrm{E})\) (V, p.~33, formula (26)).

Remark 1. --- Suppose the hilbertian space E does not reduce to 0. We identify \(\mathcal{L}(\mathbf{K};\mathbf{E})\) with E by the mapping \(u\mapsto u(1)\) ; in other words, the vector \(x\) of E is identified with the mapping \(\lambda\mapsto\lambda.x\) from K into E. Then the adjoint of \(x\) is the mapping \(x^{*}:E\to K\) given by \(x^{*}(y)=\langle x|y\rangle\) . In other words, \(x\mapsto x^{*}\) is the canonical semi-linear mapping from E onto its dual (V, p.~15).

Similarly, we identify the number \(\lambda \in \mathbf{K}\) with the endomorphism \(\lambda. 1_{\mathrm{E}}\) of E. Then \(\lambda^{*}\) is precisely the conjugate of \(\lambda\) .

With these identifications, we can define a product \(t_{1} \ldots t_{n}\) where each \(t_{i}\) is, either a number in K, or a vector in E, or a linear form belonging to \(E'\) , or an element of \(\mathcal{L}(E)\) , provided that there are never two consecutive factors \(t_{i}\) and \(t_{i+1}\) of one of the following types :

\(\bullet\) \(xy\) where \(x, y\) are both in \(\mathbf{E}\) , or both in \(\mathbf{E}'\) ;

\(\bullet\) \(x\mathbf{A}\) or \(\mathbf{A}x'\) with \(\mathbf{A} \in \mathcal{L}(\mathbf{E})\) , \(x \in \mathbf{E}\) and \(x' \in \mathbf{E}'\) .

We have the following rules of composition :

\begin{enumerate}
\def\labelenumi{\alph{enumi})}
\item
  associativity;
\item
  every element of \(\mathbf{K}\) commutes with all the other factors;
\item
  we have \((t_1 \ldots t_n)^* = t_n^* \ldots t_1^*\) ; in other words, the adjoint of a product is the product of the adjoints taken in the reverse order. Also \(t^{**} = t\) .
\end{enumerate}

For example, let \(x, y\) be in E and let A be in \(\mathcal{L}(\mathrm{E})\) . Then \(x^{*}y\) represents the scalar product \(\langle x|y\rangle\) and \(x^{*}\mathrm{A}y\) represents the scalar product \(\langle x|\mathrm{A}y\rangle\) . We also have \((\mathrm{A}^{*}x)^{*} = x^{*}\mathrm{A}^{**} = x^{*}\mathrm{A}\) , hence \((\mathrm{A}^{*}x)^{*}y = x^{*}\mathrm{A}y\) , which can be interpreted as

\[
\langle \mathrm{A} ^ {*} x | y \rangle = \langle x | \mathrm{Ay} \rangle
\]

in conformity with the definition of the adjoint. We observe that \(yx^{*}\) is the endomorphism \(z \mapsto y\langle x|z\rangle\) of E, since \(yx^{*}z\) can be interpreted as \(y(x^{*}z)\) by associativity, or as \(y.\langle x|z\rangle\) .

Following Dirac \(^{1}\) , in most works of Mathematical Physics, the elements of E are represented by the symbol \(|x\rangle\) , those of E' by \(\langle t|\) . The scalar product is written as \(\langle x|y\rangle = \langle x|.|y\rangle\) and the first rule of interdiction in the products excludes the combinations of the signs \(\rangle|\) and \(|\langle\) , for example \(|x\rangle|y\rangle\) .

\(^{1}\) See P. A. M. DIRAC, \emph{Quantum Mechanics}, Oxford University Press, New York, 1935.

PROPOSITION 3. --- Let E and F be two hilbertian spaces and \(u \in \mathcal{L}(\mathrm{E};\mathrm{F})\) . The following conditions are equivalent :

\begin{enumerate}
\def\labelenumi{(\roman{enumi})}
\item
  \(u\) is a topological vector space isomorphism, with an inverse equal to \(u^{*}\) ;
\item
  \(u\) is surjective and \(u^{*}u = 1_{\mathrm{E}}\) ;
\item
  \(u\) is injective and \(uu^{*} = 1_{\mathrm{F}}\) ;
\item
  \(u\) is an isomorphism of normed spaces;
\item
  \(u\) is a hilbertian space isomorphism.
\end{enumerate}

Condition (1) means that we have \(u^{*}u = 1_{\mathrm{E}}\) and \(uu^{*} = 1_{\mathrm{F}}\) . Hence the equivalence of (i), (ii) and (iii) follows from E, II, § 3, No.~8, prop. 8. We have already seen the equivalence of (iv) and (v) (V, p.~5). Finally, the relation \(u^{*}u = 1_{\mathrm{E}}\) is equivalent to \(\langle x|u^{*}u(y)\rangle = \langle x|y\rangle\) , that is, to \(\langle u(x)|u(y)\rangle = \langle x|y\rangle\) for all \(x,y\) in E, and evidently implies that \(u\) is injective; this proves the equivalence of (ii) and (v).

An automorphism of the hilbertian space E is also called a unitary operator, that is, an operator \(u \in \mathcal{L}(\mathrm{E})\) satisfying \(uu^{*} = u^{*}u = 1_{E}\) .

Remark 2. --- The relation \(u^{*}u = 1_{\mathrm{E}}\) does not characterize all the automorphisms of the hilbertian space E. For example, let \(\mathrm{E} = \ell^2 (\mathbf{N})\) and let \(u\) be defined by \(u(x_{n}) = x_{n - 1}\) for \(n\geqslant 1\) and \(u(x)_0 = 0\) . We have \(\| u(x)\| = \| x\|\) for all \(x\in \mathrm{E}\) , that is, \(u^{*}u = 1_{\mathrm{E}}\) , but \(u\) is not surjective.

Remark 3. --- The definition (1) of the adjoint \(u^{*}\) can also be written as

\[
\langle y | u (x) \rangle = \langle u ^ {*} (y) | x \rangle ,
\]

or, by V, p.~15, as

\[
\langle u (x), y ^ {*} \rangle = \langle x, (u ^ {*} (y)) ^ {*} \rangle .
\]

But we also have \(\langle u(x), y^* \rangle = \langle x, {}^t u(y^*) \rangle\) , hence we can express the adjoint in terms of the transpose,

\[
(u ^ {*} (y)) ^ {*} = ^ {t} u (y ^ {*}).
\]

\subsection{Partially isometric linear mappings}

DEFINITION 2. --- Let E and F be two hilbertian spaces and \(u \in \mathcal{L}(E; F)\) . The orthogonal of the kernel of u in E is said to be the initial subspace of u and the closure of the image of u in F is called the final subspace of u. The orthoprojector from E (resp. F) onto the initial (resp. final) subspace of u is called the initial (resp. final) orthoprojector of u.

Let \(\mathbf{P}\) be the initial subspace of \(u\) . Since \(\mathbf{E}\) is the direct sum of \(\mathbf{P}\) and of the kernel of \(u\) , we have \(u(\mathbf{P}) = u(\mathbf{E})\) .

PROPOSITION 4. --- (i) The initial (resp. final) subspace of \(u^*\) is equal to the final (resp. initial) subspace of \(u\) .

\begin{enumerate}
\def\labelenumi{(\roman{enumi})}
\setcounter{enumi}{1}
\tightlist
\item
  Suppose that \(\mathbf{E} = \mathbf{F}\) . Let \(\mathbf{M}\) be a closed vector subspace of \(\mathbf{E}\) and \(\mathbf{M}^{\circ}\) its orthogonal. The relations \(u(\mathbf{M}) \subset \mathbf{M}\) and \(u^{*}(\mathbf{M}^{\circ}) \subset \mathbf{M}^{\circ}\) are equivalent.
\end{enumerate}

Let \(\mathrm{Q} = \overline{u(\mathrm{E})}\) be the final subspace of \(u\) . The orthogonal \(\mathrm{Q}^{\circ}\) of \(\mathrm{Q}\) in \(\mathrm{F}\) consists of all vectors \(y\) such that \(\langle u(x)|y\rangle = 0\) for all \(x\in \mathrm{E}\) ; this is equivalent to: \(\langle x|u^{*}(y)\rangle = 0\) for all \(x\in \mathrm{E}\) , or to \(u^{*}(y) = 0\) . Hence we have \(\mathrm{Q}^{\circ} = \operatorname{Ker}u^{*}\) , and \(\mathrm{Q}\) is the initial subspace of \(u^{*}\) . Since \(u\) is the adjoint of \(u^{*}\) , the final subspace of \(u^{*}\) is also the initial subspace of \(u\) . This proves (i).

The relation \(u(\mathbf{M}) \subset \mathbf{M}\) implies that \(u(\mathbf{M})\) is orthogonal to \(\mathbf{M}^{\circ}\) , and the relation \(u^{*}(\mathbf{M}^{\circ}) \subset \mathbf{M}^{\circ}\) implies that \(u^{*}(\mathbf{M}^{\circ})\) is orthogonal to \(\mathbf{M}\) . But we have \(\langle u(x)|y\rangle = \overline{\langle u^{*}(y)|x\rangle}\) for all \(x \in \mathbf{M}\) and \(y \in \mathbf{M}^{\circ}\) ; hence (ii) follows.

We remark that prop. 4 can be deduced from the general properties of the transpose (II, p.~51, cor. 2) in view of remark 3, V, p.~41.

DEFINITION 3. --- Let E and F be two hilbertian spaces. A mapping \(u \in \mathcal{L}(\mathrm{E};\mathrm{F})\) is said to be partially isometric if \(\|u(x)\| = .\|x\|\) for all x belonging to the initial subspace of u.

Let \(u \in \mathcal{L}(\mathrm{E};\mathrm{F})\) and let \(\mathbf{N}\) be its kernel and \(\mathbf{I}\) its image. To say that \(u\) is partially isometric is the same as saying that the linear mapping \(\tilde{u}: \mathrm{E/N} \to \mathrm{I}\) deduced from \(u\) is isometric (V, p.~13). Then the subspace \(\mathbf{I}\) of \(\mathbf{F}\) is complete, hence closed, and is the final subspace of \(u\) . Consequently, \(u\) induces a hilbertian space isomorphism from the initial subspace of \(u\) onto its final subspace.

PROPOSITION 5. --- Let \(u \in \mathcal{L}(\mathrm{E};\mathrm{F})\) , let \(\mathbf{P}\) be its initial subspace and \(\mathbf{Q}\) be the final subspace. Let \(p\) (resp. \(q\) ) denote the initial (resp. final) orthoprojector of \(u\) . Assume that \(u\) is partially isometric.

\begin{enumerate}
\def\labelenumi{(\roman{enumi})}
\item
  The mapping \(u^{*} \in \mathcal{L}(\mathbf{F};\mathbf{E})\) is partially isometric, with initial subspace \(\mathbf{Q}\) and final subspace \(\mathbf{P}\) . The isomorphism from \(\mathbf{P}\) onto \(\mathbf{Q}\) induced by \(u\) is then the inverse of the isomorphism from \(\mathbf{Q}\) onto \(\mathbf{P}\) induced by \(u^{*}\) .
\item
  We have \(u^{*}u = p\) and \(uu^{*} = q\) .
\end{enumerate}

On account of prop. 4 (i), assertion (i) is a consequence of (ii).

We now prove (ii). Since \(\mathbf{P}\) contains the image of \(u^{*}\) , the mapping \(u^{*}u\) maps \(\mathbf{E}\) into \(\mathbf{P}\) . Let \(x \in \mathbf{E}\) and \(y \in \mathbf{P}\) , then

\[
\langle u ^ {*} u (x) | y \rangle = \langle u (x) | u (y) \rangle .
\]

If x belongs to P, then \(\langle u(x)|u(y)\rangle = \langle x|y\rangle\) by the definition of a partially isometric mapping; if x belongs to the kernel N of u, then \(u(x) = 0\) , hence \(\langle u(x)|u(y)\rangle = 0\) and \(\langle x|y\rangle = 0\) since N and P are orthogonal. Since \(E = P \oplus N\) , we have \(\langle u^{*}u(x) - x|y\rangle = 0\) in all the cases, and so \(u^{*}u\) is the orthoprojector p from E onto P. That \(uu^{*} = q\) follows by interchanging u and \(u^{*}\) in the above.

PROPOSITION 6. --- For every \(u \in \mathcal{L}(\mathrm{E};\mathrm{F})\) , the following conditions are equivalent :

\begin{enumerate}
\def\labelenumi{(\roman{enumi})}
\item
  \(u\) is partially isometric;
\item
  \(u^{*}\) is partially isometric;
\item
  \(u^{*}u\) is an orthoprojector;
\item
  \(uu^{*}\) is an orthoprojector;
\item
  \(uu^{*}u = u;\)
\item
  \(u^{*}uu^{*} = u^{*}\) .
\end{enumerate}

By prop. 5, (i) is equivalent to (ii).

\begin{enumerate}
\def\labelenumi{(\roman{enumi})}
\item
  \(\Rightarrow\) (v): Suppose \(u\) is partially isometric. Then \(u^{*}u\) is the initial orthoprojector of \(u\) by prop. 5. Hence for every \(x \in \mathrm{E}\) , \(u^{*}u(x) - x\) belongs to the kernel of \(u\) , that is, \(uu^{*}u(x) = u(x)\) .
\item
  \(\Rightarrow\) (iii): Suppose that \(uu^{*}u = u\) and let \(p = u^{*}u\) ; then \(p = p^{*}\) and \(p^2 = p\) . Let M (resp. N) be the image (resp. the kernel) of \(p\) . For \(x \in \mathbf{M}\) and \(y \in \mathbf{N}\) , we have \(\langle x|y\rangle = \langle p(x)|y\rangle = \langle x|p^{*}(y)\rangle = \langle x|p(y)\rangle = 0\) . Since M and N are orthogonal, \(p\) is the orthoprojector from E onto M.
\item
  \(\Rightarrow\) (i): Suppose \(p = u^{*}u\) is an orthoprojector with image M and kernel N.
\end{enumerate}

For all \(x\in \mathbf{E}\) , we have

\[
\left\| u (x) \right\| ^ {2} = \langle u ^ {*} u (x) | x \rangle = \langle p (x) | x \rangle .
\]

Hence \(u(x) = 0\) for \(x \in \mathbf{N}\) and \(\| u(x) \| = \| x \|\) for \(x \in \mathbf{M}\) , and so \(u\) is partially isometric with kernel \(\mathbf{N}\) and initial subspace \(\mathbf{M}\) .

We have proved the equivalence of (i), (iii) and (v). Replacing \(u\) by \(u^*\) , we can deduce the equivalence of (ii), (iv) and (vi). This proves prop. 6.
